\subsection{Basic notations} Given a ring $R$, we denote by $\mathcal{Z}(R)$ its center. When $R$ is commutative and has no nontrivial zero divisors, $Q(R)$ denotes its fraction field.

Given a Hopf algebra $H$ with product $m$ and coproduct $\Delta$, we denote by $H^{\rm cop}$ (resp.~$H_{\rm op}$) the Hopf algebra with the same algebra (resp.\ coalgebra) structure as $H$ but the opposite coproduct $\Delta^{\rm cop}:=\sigma \circ {\Delta}$ (resp.\ opposite product $m\circ \sigma$), where $\sigma(x\otimes y) = y\otimes x$, and the antipode ${S}^{-1}$. We use Sweedler's coproduct notation, $ \Delta(x) = \sum_{(x)}x_{(1)} \otimes x_{(2)}$, $x\in H$, and we set \smash{$\Delta^{(1)} := {\rm id}$}, \smash{$\Delta^{(2)} := \Delta$}, and \smash{$\Delta^{(n)} := (\Delta \otimes {\rm id}) \Delta^{(n-1)}$} for $n\geq 3$ (this is not the convention used in \cite{BR}).

The results of this paper hold true for any finite-dimensional complex semisimple Lie algebra~${\mathfrak g}$, but unless we state it differently, we will assume ${\mathfrak g}$ is simple. We will denote its rank by~$m$, and its Cartan matrix by~$(a_{ij})$. We fix a Cartan subalgebra $\mathfrak{h}\subset {\mathfrak g}$ and a basis of simple roots~${\alpha_i \in \mathfrak{h}_{\mr}^{*}}$, and denote by $\mathfrak{b}_\pm$ the Borel subalgebras associated to it. We denote by $N$ the number of positive roots of ${\mathfrak g}$, and by $\rho$ half the sum of the positive roots.

We denote by $d_1,\dots , d_m$ the unique coprime positive integers such that the matrix $(d_ia_{ij})$ is symmetric, and $(\ ,\ )$ the unique inner product on $\mathfrak{h}_{\mr}^{*}$ such that $d_ia_{ij} = (\alpha_i,\alpha_j)$. For any root~$\alpha$, the coroot is $\alpha\check{}=2\alpha/(\alpha,\alpha)$; in particular $\alpha\check{}_i=d_i^{-1}\alpha_i$. The root lattice $Q$ is the $\mz$-lattice in~$\mathfrak{h}_{\mr}^{*}$ defined by $ Q = \sum_{i=1}^m \mathbb{Z} \alpha_i$. The weight lattice $P$ is the $\mz$-lattice formed by all $\lambda\in \mathfrak{h}_{\mr}^{*}$ such that~${(\lambda ,\alpha\check{}_i) \in \mz}$ for every $i=1,\dots,m$. So $ P= \sum_{i=1}^m \mathbb{Z} \varpi_i$, where $\varpi_i$ is the fundamental weight dual to the simple coroot $\alpha\check{}_i$, which satisfies $( \varpi_i,\alpha\check{}_j) = \delta_{i,j}$. Note that $(\lambda,\alpha)\in \mz$ for every $\lambda\in P$, $\alpha\in Q$. We denote by $D$ the cardinality of the quotient lattice $P/Q$. Then $D$ is the smallest positive integer such that $D(\lambda,\mu)\in {\mathbb Z}$ for every $\lambda,\mu\in P$, that is, such that $DP\subset Q$.

We denote by
\[ P_+:= \sum_{i=1}^m \mathbb{Z}_{\geq 0} \varpi_i
\]
 the cone of dominant integral weights, and we put
 \[ Q_+:= \sum_{i=1}^m \mathbb{Z}_{\geq 0} \alpha_i.\]
 Though $Q\subset P$, it is not true that $Q_+\subset P_+$, but we have $DP_+ \subset Q_+$. This last property is not trivial, and follows from the classical fact that the inverse of the Cartan matrix~$(a_{ij})$ has coefficients in $D^{-1}\mn$.

We will use the standard partial order relation $\leq$ on $P$, defined by: $\lambda, \mu\in P$ satisfy $\lambda \leq \mu$ if~${\mu-\lambda \in Q_+}$. In Section \ref{HNsec}, we will also use the partial order relation $\preceq$ on $P$ defined by: $\lambda \preceq \mu$ if $\mu-\lambda \in D^{-1}Q_+$.

We denote by $\mathcal{B}(\mathfrak{g})$ the braid group of $\mathfrak{g}$; we recall its standard defining relations in Appendix~\ref{QW}.

We denote by $G$ the connected and simply-connected algebraic group with Lie algebra $\mathfrak{g}$, and by $T_G$ the maximal torus of $G$ with Lie algebra $\mathfrak{h}$; $N(T_G)$ is the normalizer of $T_G$, $W = N(T_G)/T_G$ is the Weyl group, $B_\pm$ are the Borel subgroups of $G$ with Lie algebra $\mathfrak{b}_\pm$, and $U_\pm \subset B_\pm$ are their unipotent subgroups.

We denote by $\Oo(G)$ the coordinate ring of $G$. It is a commutative Hopf algebra, which can be identified with the restricted dual of the universal enveloping algebra $U(\mathfrak{g})$ (see \cite{Ko,Lusztig3}). Similarly we denote by $\Oo(B_\pm)$ the coordinate ring of~$B_\pm$.

Let $q$ be an indeterminate, let $q^{1/D}$ be such that $\bigl(q^{1/D}\bigr)^D=q$, set $A={\mathbb C}\big[q,q^{-1}\big]$, $q_i=q^{d_i}$, $q_\beta = q^{(\beta,\beta)/2}$ for $\beta\in Q$, and given integers $p$, $k$ with $0\leq k\leq p$, we put
\begin{alignat*}{5}
 & [p]_q = \frac{q^p-q^{-p}}{q- q^{-1}} ,\qquad &&[0]_q! =1,\qquad &&[p]_q! =[1]_q[2]_q\cdots [p]_q ,\qquad &&\left[\begin{matrix} p \\ k \end{matrix}\right]_{q} =\frac{[p]_q! }{[p-k]_q![k]_q!},&\\
& (p)_q = \frac{q^p-1}{q- 1} ,\qquad&& (0)_q! =1 ,\qquad&& (p)_q! =(1)_q(2)_q\cdots (p)_q,\qquad &&\left(\begin{matrix} p \\ k \end{matrix}\right)_{q} =\frac{(p)_q! }{(p-k)_q!(k)_q!}.&
\end{alignat*}
We denote by $\mathcal{A}_0\subset \mc(q)$ the ring of functions regular at $q=0$; this ring is used only in Section~\ref{canbasemodqg}.

We denote by $\e$ a primitive $l$-th root of unity such that $\e^{2d_i}\ne 1$ is also a primitive $l$-th root of unity for all $i\in \{1,\dots,m\}$. This means that $l$ is odd, and coprime to $3$ if $\mathfrak{g}$ is $G_2$. We put $\e_i:=\e^{d_i}$.

In this paper, we use the definition of the unrestricted integral form $U_A(\mathfrak{g})$ given in \cite{DC-L,DC-K-P2}; in \cite{BR} we used the one of \cite{DC-K,DC-K-P1}. The two are (trivially) isomorphic, and have the same specialization at $q=\e$. Also, we denote here by $L_i$ the generators of $U_q(\mathfrak{g})$ we denoted by $\ell_i$ in~\cite{BR}.

In order to facilitate the comparison with the results of~\cite{DC-L}, we note that their generators denoted $K_i$, $E_i$ and $F_i$, that we will denote by $\tilde{K}_i$, $\tilde{E}_i$ and $\tilde{F}_i$, can be written as $K_i$, $K_i^{-1}E_i$ and~$F_iK_i$ in our notations. They satisfy the same algebra relations.

\section{Background results}\label{PREL}

\subsection[On U\_q, O\_q, L\_{0,n}, M\_{0,n}, and Phi\_n]{On $\boldsymbol{U_q}$, $\boldsymbol{\Oo_q}$, $\boldsymbol{\Ll_{0,n}}$, $\boldsymbol{\Mm_{0,n}}$, and $\boldsymbol{\Phi_n}$}\label{ALGdef}
Except when stated differently, we refer to \cite[Sections 2--4 and 6]{BR}, and the references therein for details about the material of this section. We stress that the simply-connected quantum group, that we denote $U_q$ below, was denoted $\tilde U_q$ in \cite{BR}. Also, the adjoint quantum group $U_q^{\rm ad}$ was denoted $U_q$.


The {\it simply-connected} quantum group $U_q = U_q(\mathfrak{g})$ is the Hopf algebra over $\mathbb {C}(q)$ with generators $E_i$, $F_i$, $L_i$, $L_i^{-1}$, $1\leq i \leq m$, and defining relations
\begin{gather*}%\label{EFK}\label{Serre1}\label{Serre2}
L_iL_j=L_jL_i , \qquad L_iL_i^{-1}=L_i^{-1}L_i=1,\qquad L_iE_jL_i^{-1}=q_i^{\delta_{i,j}}E_j ,\qquad L_iF_jL_i^{-1}=q_i^{-\delta_{i,j}}F_j, \\
E_iF_j-F_jE_i=\delta_{i,j}\frac{K_i-K_i^{-1}}{q_i-q_i^{-1}},\\
\sum_{r=0}^{1-a_{ij}} (-1)^r \left[\begin{matrix} 1-a_{ij} \\ r \end{matrix}\right]_{q_i} E_i^{1-a_{ij}-r}E_jE_i^{r} = 0 \qquad {\rm if}\quad i\ne j,\\
\sum_{r=0}^{1-a_{ij}} (-1)^r \left[\begin{matrix} 1-a_{ij} \\ r \end{matrix}\right]_{q_i} F_i^{1-a_{ij}-r}F_jF_i^{r} = 0 \qquad {\rm if}\quad i\ne j,
\end{gather*}
where for $ \lambda=\sum_{i=1}^m m_i \varpi_i \in P$ we set $ K_\lambda=\prod_{i=1}^m L_i^{m_i}$, and $ K_i=K_{\alpha_i}=\prod_{j=1}^m L_j^{a_{ji}}$.
The coproduct $\Delta$, antipode $S$, and counit $\varepsilon$ of $U_q$ are given by
\begin{gather*}
\Delta(L_i)=L_i\otimes L_i ,\qquad \Delta(E_i)=E_i\otimes K_i+1\otimes E_i ,\qquad \Delta(F_i)=F_i\otimes 1 + K_i^{-1}\otimes F_i, \\ S(E_i) = -E_iK_i^{-1},\qquad S(F_i) = -K_iF_i ,\qquad S(L_i) = L_i^{- 1},\\ \varepsilon(E_i) = \varepsilon(F_i)=0,\qquad \varepsilon(L_i)=1.
\end{gather*}
We fix a reduced expression $s_{i_1}\cdots s_{i_N}$ of the longest element $w_0$ of the Weyl group of $\mathfrak{g}$. It induces a total ordering of the positive roots,
\[\beta_1 = \alpha_{i_1},\qquad \beta_2 = s_{i_1}(\alpha_{i_2}), \qquad \dots, \qquad\beta_N = s_{i_1}\cdots s_{i_{N-1}}(\alpha_{i_N}).\]
The root vectors of $U_q$ with respect to such an ordering are defined by
\begin{equation}\label{rootvectdef}
E_{\beta_k} = T_{i_1}\cdots T_{i_{k-1}}(E_{i_k}) ,\qquad F_{\beta_k} = T_{i_1}\cdots T_{i_{k-1}}(F_{i_k}),
\end{equation}
where $T_i$ is the Lusztig algebra automorphism of $U_q$ associated to the simple root $\alpha_i$ \mbox{\cite{Lusztig2,Lusztig}} (see also \cite[Chapter~8]{CP}). The braid group $\mathcal{B}(\mathfrak{g})$ acts on $U_q$ by means of the Lusztig automorphisms. In the appendix, we recall the relation between $T_i$ and the generator $\hat{w}_i$ of the quantum Weyl group, which we will mostly use. Let us just recall here that the monomials $F_{\beta_1}^{r_1}\cdots F_{\beta_N}^{r_N}K_\lambda E_{\beta_N}^{t_N}\cdots E_{\beta_1}^{t_1}$ ($r_i,t_i\in \mn$, $\lambda\in P$) form a basis of $U_q$, the {\it PBW basis}.

$U_q$ is a {\it pivotal} Hopf algebra, with pivotal element \smash{$ \ell :=K_{2\rho}=\prod_{j=1}^m L_j^2$}.
So $\ell$ is group-like, and $S^2(x) = \ell x\ell^{-1}$ for every $x\in U_q$.

The {\it adjoint} quantum group $U_q^{\rm ad} = U_q^{\rm ad}(\mathfrak{g})$ is the Hopf subalgebra of $U_q$ generated by the elements $E_i$, $F_i$ ($i=1,\dots,m$) and $K_\alpha$ with $\alpha\in Q$; so $\ell\in U_q^{\rm ad}$. When $\mathfrak{g}={\mathfrak{sl}_2}$, we simply write the above generators $E=E_1$, $F=F_1$, $L=L_1$, $K=K_1$.

We denote by $U_q(\mathfrak{n}_+)$, $U_q(\mathfrak{n}_-)$ and $U_q(\mathfrak{h})$ the subalgebras of $U_q$ generated respectively by the~$E_i$, the $F_i$, and the $K_\lambda$ ($\lambda\in P$), and by $U_q(\mathfrak{b}_+)$ and $U_q(\mathfrak{b}_-)$ the subalgebras generated by the~$E_i$ and the $K_\lambda$, and by the $F_i$ and the $K_\lambda$, respectively. We do similarly with $U_q^{\rm ad}$, where now $U_q^{\rm ad}(\mathfrak{h})$ is generated by the $K_\lambda$ with~$\lambda\in Q$.

The Hopf algebra $U_q^{\rm ad}$ is not braided in a strict sense, but it has braided categorical completions. Let us recall briefly what this means and implies. For details, we refer to \cite[Sections~2 and~3]{BR} (see also \cite[Section~3.10]{VY}, where $\mUq$ below is formulated in terms of multiplier Hopf algebras).

A $U_q^{\rm ad}$-module $V$ is said {\it of type $1$} if it has finite dimension and the generators $K_i$ are diagonalizable on $V$ with eigenvalues in \smash{$q_i^{\mz}$}. We denote by $\mathcal{C}$ the category of $U_q^{\rm ad}$-modules of type $1$, by~${\rm Vect}$ the category of finite-dimensional $\mc(q)$-vector spaces, and by $F_{\mathcal C}\colon{\mathcal C}\to {\rm Vect}$ the forgetful functor. The category $\mathcal{C}$ is semisimple. The simple objects are highest weight $U_q^{\rm ad}$-modules; we denote by $V_\mu$ the simple module with highest weight $\mu\in P_+$. In the case $\mathfrak{g}={\mathfrak{sl}_2}$, we identify~$P_+$ with $\mathbb{N}$, and denote by $V_n$ the simple module of dimension $n+1$. Note that $V_\mu$ is canonically endowed with a structure of $U_q$-module of type $1$, the generators $L_i$ being diagonalizable with eigenvalues in \smash{$q_i^{\mz/D}$}. The {\it categorical completion} $\mUq^{\rm ad}$ of $U_q^{\rm ad}$ is the set of natural transformations~${F_{\mathcal C}\ra F_{\mathcal C}}$. An element of \smash{$\mUq^{\rm ad}$} is a collection $(a_V)_{V\in {\rm Ob}(\mathcal{C})}$, where $a_V\in {\rm End}_{\mc(q)}(V)$ satisfies~${F_{\mathcal C}(f)\circ a_V=a_W\circ F_{\mathcal C}(f)}$ for any objects $V$, $W$ of $\mathcal{C}$ and any arrow \smash{$f\in {\rm Hom}_{U_q^{\rm ad}}(V,W)$}. It is not hard to see that $\mUq^{\rm ad}$ inherits from $\mathcal{C}$ a natural structure of (completed) Hopf algebra such that the map
\begin{gather}\label{defiota}
\iota\colon\ U_q^{\rm ad} \longrightarrow \mUq^{\rm ad}, \qquad x\longmapsto(\pi_V(x))_{V\in {\rm Ob}(\mathcal{C})}
\end{gather}
is a morphism of Hopf algebras, where $\pi_V\colon U_q^{\rm ad} \ra {\rm End}(V)$ is the representation associated to a module $V$ in $\mathcal{C}$. It is a theorem that this map is injective. From now on, let us extend the coefficient ring of the modules and morphisms in $\mathcal{C}$ to $\mc\big(q^{1/D}\big)$. Put \smash{$\mUq = \mUq^{\rm ad} \bigotimes_{\mc(q)} \mc\big(q^{1/D}\big)$}.
The map $\iota$ above extends to an embedding of $U_q$ in $\mUq$. The category $\mathcal{C}$, with coefficients extended to $\mc\big(q^{1/D}\big)$, is braided and ribbon; we postpone a discussion of that fact to Section \ref{intdual}, where it will be developed. As a consequence, we can regard $\mUq$ as a quasitriangular and ribbon Hopf algebra in a generalized sense (see \cite{BR}). The $R$-matrix of $\mUq$ is the family of morphisms
\[R = (R_{V,W})_{V,W\in {\rm Ob}(\mathcal{C})},\]
where $R_{V,W}\in {\rm End}(V\otimes W)$ is the endomorphism defined by the action of Drinfeld's universal $R$-matrix on $V\otimes W$. The ribbon element of $\mUq$ is defined similarly by Drinfeld's universal ribbon element. One defines the {\it categorical tensor product} \smash{$\mUq^{\hat{\otimes} 2}$} similarly as $\mUq$; in particular it contains all the infinite series of elements of $\mUq^{\otimes 2}$ having only a finite number of non-zero terms when evaluated on a given module $V\otimes W$ of $\mathcal{C}$. There is an expansion of $R$ as an infinite series in~\smash{$\mUq^{\hat{\otimes} 2}$}. Adapting Sweedler's coproduct notation $ \Delta(x)=\sum_{(x)}x_{(1)}\otimes x_{(2)}$, we find convenient to write this series as
\begin{equation}\label{Rcat}
R=\sum_{(R)}R_{(1)}\otimes R_{(2)}.
\end{equation}
We put $R^+:=R$, $R^- := (\sigma\circ R)^{-1}$ where $\sigma$ is the flip map $x\otimes y \mapsto y\otimes x$. We will not use any explicit formula of $R$, but the following factorization formula
\begin{equation}\label{Rmatfact}
R=\Theta \hat{R},
\end{equation}
where
\[\Theta = q^{\sum_{i,j=1}^m (B^{-1})_{ij} H_i\otimes H_j} \in \mUq^{\hat{\otimes} 2},\]
with $B\in M_m(\mathbb{Q})$ the matrix with entries $B_{ij}:=d_j^{-1}a_{ij}$, and
\[\hat{R}=\sum_{(\hat{R})} \hat{R}_{(1)} \otimes \hat{R}_{(2)}\in \mUq(\mathfrak{n}_+)\hat{\otimes} \mUq(\mathfrak{n}_-)\]
(see \cite[Section 3.2]{BR}, and for details, e.g., \cite[Theorem 8.3.9]{CP}, or \cite[Theorem 3.108]{VY}). If $x$,~$y$ are weight vectors of weights $\mu$, $\nu$ respectively, then $\Theta(x\otimes y) = q^{(\mu,\nu)} x\otimes y$. Moreover, $\hat{R}$ has weight $0$ for the adjoint action of $U_q(\mathfrak{h})$; that is, complementary components $\hat{R}_{(1)}$ and $\hat{R}_{(2)}$ have opposite weights.


Recall that we denote by $G$ the connected and simply-connected algebraic group with Lie algebra $\mathfrak{g}$. The {\it quantum function Hopf algebra} $\Oo_q=\Oo_q(G)$ is defined as the restricted dual of~$U_q^{\rm ad}$ with respect to the category $\mathcal{C}$, that is, the set of ${\mathbb C}(q)$-linear maps $f\colon U_q^{\rm ad}\ra \mc(q)$ such that ${\rm Ker}(f)$ contains a cofinite two sided ideal $I$ (i.e., such that $I\oplus M =U_q$ for some finite-dimensional vector space $M$), and $ \prod_{s=-r}^r (K_i - q_i^s) \in I$ for some $r\in \mathbb{N}$ and every $i$ (see, e.g., \cite[Chapter I.7]{BG}).

The space $\Oo_q$ is a Hopf algebra, with structure maps defined dually to those of $U_q^{\rm ad}$. We denote by $\star$ its product. The algebras $\Oo_q(T_G)$, $\Oo_q(U_\pm)$, $\Oo_q(B_\pm)$ are defined similarly, by replacing~$U_q^{\rm ad}$ with $U_q^{\rm ad}(\mathfrak{h})$, $U_q^{\rm ad}(\mathfrak{n}_\pm)$, $U_q^{\rm ad}(\mathfrak{b}_\pm)$, respectively. As a vector space, $\Oq$ is generated by the functionals $x\mapsto w(\pi_V(x)v)$, $x\in U_q^{\rm ad}$, for every object $V\in {\rm Ob}(\mathcal{C})$ and vectors $v\in V$, $w\in V^*$. Such functionals are called {\it matrix coefficients}. Because the morphism $\iota\colon U_q^{\rm ad}\to \mathbb{U}_q$ is injective (see~\eqref{defiota}), the Hopf duality pairing $\langle\cdot,\cdot\rangle \colon \Oq \times U_q^{\rm ad} \ra \mc(q)$ is non degenerate. By extending the coefficient ring from $\mc(q)$ to $\mc\big(q^{1/D}\big)$, we can uniquely extend it to a bilinear pairing
\[\langle\cdot,\cdot\rangle\colon\ \big(\Oq \bigotimes_{{\mathbb C}(q)} {\mathbb C}\big(q^{1/D}\big)\big)\times \mUq \ra \mc\big(q^{1/D}\big)\]
such that the following diagram is commutative:


\centerline{\xymatrix{
 \Oq \otimes U_q^{\rm ad} \ar[r]^{\langle\cdot,\cdot \rangle}\ar[d]_{{\rm id} \otimes \iota} & {\mathbb C}(q) \ar[d]\\
 \bigl(\Oq \bigotimes_{{\mathbb C}(q)} {\mathbb C}\big(q^{1/D}\big) \bigr)\otimes \mUq \ar[r]^{\;\;\;\;\;\;\;\ \;\ \;\;\;\ \langle\cdot,\cdot \rangle}& {\mathbb C}\big(q^{1/D}\big).
 }}

This pairing is defined by $\langle {}_Y\phi{}^w_v , (a_{X})\rangle=w(a_{Y}v)$
for every $(a_{X})\in \mUq$ and ${}_Y\phi{}^w_v\in \Oq$. It is non degenerate.


The maps
\begin{gather}\label{phipm}
\Phi^\pm\colon\ \Oo_q \longrightarrow U_q^{\rm cop},\qquad \alpha\longmapsto (\alpha \otimes {\rm id})\big(R^\pm\big)= \sum_{(R^\pm)} \big\langle \alpha , R_{(1)}^\pm\big\rangle R_{(2)}^\pm
\end{gather}
are well-defined morphisms of Hopf algebras. Here we stress that it is the simply-connected quantum group $U_q^{\rm cop}$ that is the range of $\Phi^\pm$. This will be explained with more details in Section~\ref{intdual}.

Let us make two simple observations, for future reference. Firstly, because $\Oo_q$ is spanned by the matrix coefficients of the objects of $\mathcal{C}$, and $\mathcal{C}$ is semisimple with simple objects the $U_q^{\rm ad}$-modules $V_\mu$, $\mu\in P_+$, there is a decomposition of $U_q$-bimodule
\begin{equation}\label{decompdirectOq}
 \Oo_q = \bigoplus_{\mu\in P_+} C(\mu),
 \end{equation}
 where $C(\mu) = V_\mu^* \otimes V_\mu$, the space of matrix coefficients of $V_\mu$, is endowed with the left action on the factor $V_\mu$ and the right action on $V_\mu^*$, and $\Oo_q$ has the left and right coregular actions $\lhd$ and~$\rhd$, defined by
\[x\rhd \alpha := \sum_{(\alpha)}\alpha_{(1)} \langle \alpha_{(2)},x\rangle,\qquad \alpha \lhd x := \sum_{(\alpha)} \langle \alpha_{(1)},x\rangle \alpha_{(2)}\]
for all $x\in U_q$ and $\alpha\in \Oq$. Here we recall that each $U_q^{\rm ad}$-module $V_\mu$ can be regarded as a~$U_q$-module, so the above expressions make sense. The decomposition \eqref{decompdirectOq} is the {\it Peter--Weyl} decomposition of $\Oo_q$. It will be refined in Section \ref{canbasemodqg}.

Moreover, the algebra $\Oo_q$ is generated by the matrix coefficients of the simple $U_q^{\rm ad}$-modules~$V_{\!\varpi_{\!k}}$ with highest weights the fundamental weights $\varpi_k$, $k=1,\dots,m$; see, e.g., \cite[Proposition~I.7.8]{BG} for a proof. This relies on the standard fact that, for any $\mu,\nu \in P_+$ we have a direct sum decomposition of modules (where $m(\lambda)\in \mn$)
\begin{equation}\label{decomprep}
V_\mu\otimes V_\nu = V_{\mu+\nu} \oplus \bigoplus_{\lambda < \mu+\nu} V_\lambda^{\oplus m(\lambda)}.
\end{equation}
In particular, this decomposition implies that, up to scalar multiples, there is a unique non-zero morphism $V_{\mu+\nu} \ra V_\mu\otimes V_\nu$, which is injective and splits. Dually, this means that, applying the product in $\Oo_q$ followed by the projection onto the subspace $C(\mu+\nu)$ we get a canonical projection map
\begin{equation}\label{projOq}
C(\mu)\otimes C(\nu) \ra C(\mu+\nu).
\end{equation}
The {\it loop algebra} $\Ll_{0,1} = \Ll_{0,1}(\mathfrak{g})$ is defined by twisting the product $\star$ of $\Oo_q$, keeping the same underlying linear space. The new product is equivariant with respect to the right coadjoint action ${\rm coad}^r$ of $U_q$, defined by
\[{\rm coad}^r(x)(\alpha) = \sum_{(x)}{S}(x_{(2)}) \rhd \alpha \lhd x_{(1)}\]
for all $x\in U_q$ and $\alpha\in \Oq$. By equivariant we mean that $\Ll_{0,1}$ is a $U_q$-module algebra. Let us spell out its product and equivariance property. Using the fact that $U_q$ can be regarded as a~subspace of $\mUq$, the actions $\lhd$ and $\rhd$ extend naturally to actions of $\mUq$, and the product of~$\Ll_{0,1}$ is expressed in terms of $\star$ by the formula (see \cite[Proposition 4.1]{BR}):
\begin{equation}\label{mmtilde}
\alpha\beta = \sum_{(R),(R)}(R_{(2')}{S}(R_{(2)}) \rhd \alpha) \star (R_{(1')}\rhd \beta \lhd R_{(1)}),
\end{equation}
where $ \sum_{(R)}R_{(1) }\otimes R_{(2)}$ and $\sum_{(R)}R_{(1') }\otimes R_{(2')}$ are expansions of two copies of $R\in \mUq^{\hat \otimes 2}$. Note that the sum in \eqref{mmtilde} has only a finite number of non-zero terms. By using that
\[R\Delta = \Delta^{\rm cop}R,\]
 this product can equivalently be expressed as
\begin{equation}\label{mmtildebis}
\alpha\beta = \sum_{(R),(R)}(\beta \lhd R_{(1)}R_{(1')}) \star (S(R_{(2)})\rhd \alpha \lhd R_{(2')}).
\end{equation}
This product gives $\Ll_{0,1}$ (like $\Oo_q$) a structure of $U_q$-module algebra for the actions $\rhd$, $\lhd$, but also for ${\rm coad}^r$ (which is not the case of $\Oo_q$). Spelling this out for ${\rm coad}^r$, this means
\[{\rm coad}^r(x)(\alpha\beta)=\sum_{(x)}{\rm coad}^r(x_{(1)})(\alpha){\rm coad}^r(x_{(2)})(\beta).\]
The relations between $\Oo_q$, $\Ll_{0,1}$ and $U_q$ are encoded by the map
\begin{equation}\label{phi1def}
\Phi_1\colon\ \Oo_q\longrightarrow \mUq, \qquad \alpha\longmapsto(\alpha \otimes {\rm id})(RR'),
\end{equation}
where $R' = \sigma\circ R$, and as usual $\sigma\colon x\otimes y \mapsto y\otimes x$. Note that
\begin{equation}\label{RSDmapdef}
\Phi_1 = m\circ \big(\Phi^+\otimes \big(S^{-1}\circ \Phi^-{}\big)\big)\circ \Delta.
\end{equation}
We call $\Phi_1$ the {\it RSD} map, for Drinfeld, Reshetikhin and Semenov-Tian-Shansky introduced it first (see \cite{Dr,Maj,RSTS}). It is a fundamental result of the theory (see \cite{Bau1,Ca,Jos}) that $\Phi_1$ affords an isomorphism of $U_q$-modules
$\Phi_1\colon \Oo_q \ra U_q^{\rm lf}$.
For full details on that result we refer to \cite[Section 3.12]{VY}. Here, $U_q^{\rm lf}$ is the set of \emph{locally finite} elements of $U_q$, endowed with the right adjoint action ${\rm ad}^r$ of $U_q$. It is defined by
\[U_q^{\rm lf} :=\{x\in U_q\mid {\rm rk}_{\mc(q)}({\rm ad}^r(U_q)(x)) < \infty\}\]
and
\[{\rm ad}^r(y)(x) = \sum_{(y)}{S}(y_{(1)}) x y_{(2)}\]
for every $x,y \in U_q$. The action ${\rm ad}^r$ gives in fact $U_q^{\rm lf}$ a structure of right $U_q$-module algebra. It is also a right coideal, that is \smash{$\Delta\big(U_q^{\rm lf}\big) \subset U_q^{\rm lf}\otimes U_q$}. Moreover, $\Phi_1$ affords an isomorphism of $U_q$-module algebras \smash{$
\Phi_1\colon \Ll_{0,1} \ra U_q^{\rm lf}$}.
It is a fact that $\Phi_1$ affords an isomorphism between the centers $\mathcal{Z}(\Ll_{0,1})$ of $\Ll_{0,1}$ and $\mathcal{Z}(U_q)$ of $U_q$ \cite[Proposition~6.24]{BR}. Since $\Phi_1$ is an isomorphism of $U_q$-modules and \smash{$\mathcal{Z}(U_q)=U_q^{U_q}$}, it follows that \smash{$\mathcal{Z}(\Ll_{0,1}) = \Ll_{0,1}^{U_q}$}.

Let us recall a few fundamental results about $U_q^{\rm lf}$ that we will meet again later. Denote by~${T\subset U_q}$ the multiplicative Abelian group formed by the elements $K_{\lambda}$, $\lambda \in P$, and by $T_2\subset T$ the subgroup formed by the elements $K_{\lambda}$, $\lambda \in 2P$. Consider the subset $T_{2-} \subset T_2$ formed by the elements $K_{-\lambda}$, $\lambda \in 2P_+$. Clearly, $T_2 = T_{2-}^{-1}T_{2-}$ and ${\rm Card}(T/T_2) = 2^m$.

\begin{teo} \label{JLteo1} \quad
\begin{itemize}\itemsep=0pt
\item[$(1)$] $U_q^{\rm lf} = \bigoplus_{\lambda \in 2P_+} {\rm ad}^r(U_q)(K_{-\lambda})$.
\item[$(2)$] $U_q = T_{2-}^{-1}U_q^{\rm lf}[T/T_{2}]$, so $U_q$ is a free $T_{2-}^{-1}U_q^{\rm lf}$-module of rank $2^m$.
\item[$(3)$] The ring $U_q^{\rm lf}$ is $($left and right$)$ Noetherian.
\end{itemize}
\end{teo}

These results were proved by Joseph--Letzter in \cite[Theorem~4.10]{JL2}, \cite[Theorem 6.4]{JL}, and \cite[Theorem~7.4.8]{Jos}, respectively (see also \cite[Sections~7.1.6, 7.1.13 and~7.1.25]{Jos}). For~(1) and~(3), we refer also to \cite[Theorems~3.113 and~3.137]{VY}, which provides simpler proofs. For instance, in the ${\mathfrak{sl}_2}$ case, we have
\[U_q({\mathfrak{sl}_2}) = U_q({\mathfrak{sl}_2})^{\rm lf}[K] \oplus U_q({\mathfrak{sl}_2})^{\rm lf}[K].L.\]
The actual values of $\Phi_1$ are complicated in general, however, there is a simple important one, that we describe now. Let $V_{-\lambda}$ be the type $1$ simple $U_q^{\rm ad}$-module of lowest weight $-\lambda\in -P_+$ (i.e., the highest weight \smash{$U_q^{\rm ad}$}-module $V_{-w_0(\lambda)}$ of highest weight $-w_0(\lambda)$, where $w_0$ is the longest element of the Weyl group; note that $-w_0$ permutes the simple roots). Let $v\in V_{-\lambda}$ be a lowest weight vector, and $v^*\in V_{-\lambda}^*$ be such that $v^*(v)=1$ and $v^*$ vanishes on a $U_q^{\rm ad}(\mathfrak{h})$-invariant complement of $v$. Define $\psi_{-\lambda}\in \Oo_q$ by $\langle \psi_{-\lambda},x\rangle = v^*(xv)$, $x\in U_q$. From the definition~\eqref{phi1def}, it is quite easy to see that
\begin{equation}\label{Phi1psi-}
\Phi_1(\psi_{-\lambda}) = K_{-2\lambda}.
\end{equation}
In particular, $\Phi_1(\psi_{-\rho}) = \ell^{-1}$, where as usual $\ell$ is the pivotal element of $U_q$.

\begin{Remark}\label{remsocle} Since $\Ll_{0,1}=\Oo_q$ as a vector space, we still denote by $C(\mu)$, $\mu\in P^+$, the linear subspace generated by the matrix coefficients of $V_\mu$, the $U_q^{\rm ad}$-module of type $1$ and highest weight~$\mu$. It can be proved (see \cite[Section~7.1.22]{Jos}, or \cite[p.~156]{VY}, where different conventions are used) that $\Phi_1$ yields an isomorphism of $U_q$-modules
\begin{equation}\label{restphi1}
\Phi_{1}\colon\ C(-w_0(\mu)) \ra {\rm ad}^r(U_q)(K_{-2\mu}).
\end{equation}
Therefore, the summands in (1) are finite-dimensional \smash{$U_q$}-modules, and the action ${\rm ad}^r$ is completely reducible on \smash{$U_q^{\rm lf}$}. In fact, \smash{$U_q^{\rm lf}$} is the socle of ${\rm ad}^r$ on $U_q$.
\end{Remark}

\begin{Remark}\label{factell} Because $\ell =\prod_{j=1}^m L_j^2$ and $\Phi_1(\psi_{-\rho}) = \ell^{-1}$, a natural question is the factorization of~$\psi_{-\rho}$ in $\Ll_{0,1}$ (see Corollary \ref{Phi1int2}). This question is considered in \cite{JW}, where ${\mathcal L}_{0,1}({\mathfrak g})$ for~${{\mathfrak g}=\mathfrak{gl}(r+1)}$ is analysed and quantum minors are extensively studied. Let us review here some of their results in relation with $\psi_{-\rho}$.

First note that for ${\mathfrak g}=\mathfrak{sl}(r+1)$ the irreducible representation $V_{-\rho}$ of lowest weight~${-\rho}$ is isomorphic to the representation of highest weight $V_{\rho}$ because $-w_0(\rho)=\rho$. By the Weyl formula, the dimension of this representation is
\[ \prod_{\alpha>0}\frac{(2\rho,\alpha)}{(\rho,\alpha)}=2^N.
\]
 In \cite{KS}, a presentation of~${U_q(\mathfrak{gl}(r\!+1))}$ is given, which differs from our presentation of ${U_q(\mathfrak{sl}(r\!+1))}$ only by its subalgebra~${U_q({\mathfrak h})}$, generated by $r+1$ elements ${\mathbb K}_1,\dots , {\mathbb K}_{r+1}$. The inclusion
 \[U_q(\mathfrak{sl}(r+1))\subset U_q(\mathfrak{gl}(r+1))
 \]
 is such that~${K_i={\mathbb K}_i^2{\mathbb K}_{i+1}^{-2}}$, $i=1,\dots,r$. The quantum minors, properly defined in \cite{JW}, of the matrix of matrix elements of the natural representation of $U_q(\mathfrak{gl}(r+1))$ are denoted $\det_q(A_{\geq k})$ for~${k=1,\dots ,r+1}$. We have $\det_q(A_{\geq 1})=1$ in the case of $\mathfrak{sl}(r+1)$. Then \cite{JW} proves that $\det_q(A_{\geq k})=({\mathbb K}_k\cdots{\mathbb K}_{r+1})^2$, and there exists an element $\mathbb{K}\in U_q(\mathfrak{gl}(r+1))$ such that
\[{\mathbb K}^{-2\rho}={\det}_q(A_{\geq 1})^{-r}{\det}_q(A_{\geq 2})\cdots {\det}_q(A_{\geq r+1}).\]
This has to be interpreted as $K_{-2\rho}=\Phi_1(\det_q(A_{\geq 2})\cdots \det_q(A_{\geq r+1}))$ in the case of $\mathfrak{sl}(r+1)$. As a result, this gives the equality
\[
\psi_{-\rho}={\det}_q(A_{\geq 2})\cdots {\det}_q(A_{\geq r+1}).
\]
\end{Remark}

The {\it $($quantum$)$ graph algebra} $\Ll_{0,n} = \Ll_{0,n}(\mathfrak{g})$ is the braided tensor product of $n$ copies of~$\Ll_{0,1}$ (considered as a $U_q$-module algebra). As a linear space and $U_q$-bimodule with actions $\lhd$ and~$\rhd$, it coincides with $\Ll_{0,1}^{\otimes n}$, and thus with $\Oo_q^{\otimes n}$. It is also a right $U_q$-module algebra, with the following action of $U_q$ (extending ${\rm coad}^r$ on $\Ll_{0,1}$):
 \begin{gather}
{\rm coad}_n^r(y)\big(\alpha^{(1)}\otimes \dots \otimes \alpha^{(n)}\big) =
\sum_{(y)} {\rm coad}^r(y_{(1)})\big(\alpha^{(1)}\big) \otimes \dots \otimes {\rm coad}^r(y_{(n)})\big(\alpha^{(n)}\big)\label{actionprod}
\end{gather}
for all $y \in U_q$ and $\alpha^{(1)}\otimes \dots \otimes \alpha^{(n)} \in \Ll_{0,n}$. The product of $\Ll_{0,n}$ can be expressed as follows. For every $1\leq a\leq n$, define $\mathfrak{i}_a\colon \Ll_{0,1}\ra \Ll_{0,n}$ by $\mathfrak{i}_a(x) = 1^{\otimes (a-1)}\otimes x \otimes 1^{\otimes (n-a)}$; $\mathfrak{i}_a$ is an embedding of~$U_q$-module algebras. We will use the notations
\begin{equation}\label{plgtL01a}
\Ll_{0,n}^{(a)}:= {\rm Im}(\mathfrak{i}_a),\qquad (\alpha)^{(a)} := \mathfrak{i}_a(\alpha).
\end{equation}
Take $(\alpha)^{(a)},(\alpha')^{(a)} \in \Ll_{0,n}^{(a)}$ and $(\beta)^{(b)},(\beta')^{(b)} \in \Ll_{0,n}^{(b)}$ with $a<b$. Then the product of $\Ll_{0,n}$ is given by the following formula (see \cite[Section 6]{BR}):
\begin{align}
\big((\alpha)^{(a)}\otimes (\beta)^{(b)}\big) \big((\alpha')^{(a)} \otimes (\beta')^{(b)}\big) ={}& \sum_{(R^1),\dots, (R^4)} \big(\alpha \big(S\big(R^3_{(1)}R^4_{(1)}\big)\rhd \alpha' \lhd R^1_{(1)}R^2_{(1)}\big)\big)^{(a)}\nonumber\\
 &
\otimes \big( \big( S\big(R^1_{(2)}R^3_{(2)}\big) \rhd \beta \lhd R^2_{(2)}R^4_{(2)}\big)\beta'\big)^{(b)},\label{prodL0nform}
\end{align}
where $R^i = \sum_{(R^i)}R_{(1)}^i\otimes R_{(2)}^i$, $i\in \{1,2,3,4\}$, are expansions of four copies of $R\in \mUq^{\hat \otimes 2}$, and on the right-hand side the product is componentwise that of $\Ll_{0,1}$. Later we will use the fact that the product of $\Ll_{0,n}$ is obtained from the standard (componentwise) product of $\Ll_{0,1}^{\otimes n}$ by a process that may be inverted. Indeed, \eqref{prodL0nform} can be rewritten as
\begin{equation}\label{prodL0nform2}
\big((\alpha)^{(a)}\otimes (\beta)^{(b)}\big) \big((\alpha')^{(a)} \otimes (\beta')^{(b)}\big) = \sum_{(F)} (\alpha)^{(a)}\big((\alpha')^{(a)}\cdot F_{(2)}\big)\otimes \big((\beta)^{(b)}\cdot F_{(1)}\big)(\beta')^{(b)},
\end{equation}
where $F=\sum_{(F)} F_{(1)}\otimes F_{(2)} := (\Delta \otimes \Delta)(R')$, and the symbol ``$\cdot$'' stands for the right action of \smash{$\mUq^{\hat \otimes 2}$} on $\Ll_{0,1}$ that may be read from \eqref{prodL0nform}. The tensor $F$ is known as a twist. Then, by replacing $F$ with its inverse $\bar{F} = (\Delta \otimes \Delta)\big(R'{}^{-1}\big)$, one can express the product of
$\Ll_{0,1}^{\otimes n}$ in terms of the product of $\Ll_{0,n}$ by
\begin{equation}\label{prodL0nform3}
(\alpha)^{(a)}(\alpha')^{(a)}\otimes (\beta)^{(b)}(\beta')^{(b)} = \sum_{(\bar{F})}\big((\alpha)^{(a)}\otimes
 \big((\beta)^{(b)}\cdot \bar{F}_{(1)}\big)\big) \big( \big((\alpha')^{(a)}\cdot \bar{F}_{(2)}\big)\otimes (\beta')^{(b)}\big).
\end{equation}
 We call {\it quantum moduli algebra} and denote by $\Mm_{0,n} = {\mathcal L}_{0,n}^{U_q}$ the subalgebra of $\Ll_{0,n}$ formed by the $U_q$-invariant elements.


The map $\Phi_1$ can be extended to $\Ll_{0,n}$ as follows. Consider the following action of $U_q$ on the tensor product algebra $U_q^{{\otimes} n}$, which extends ${\rm ad}^r$ on $U_q$:
\[% \label{adjointactiononn}
{\rm ad}_n^r(y)(x) = \sum_{(y)}\Delta^{(n)}({S}(y_{(1)})) x \Delta^{(n)}(y_{(2)})
\]
for all $y \in U_q$, $x \in U_q^{{\otimes} n}$. This action gives $U_q^{{\otimes} n}$ a structure of right $U_q$-module algebra. In~\cite{A}, Alekseev introduced a morphism of $U_q$-module algebras $\Phi_n\colon \Ll_{0,n} \ra U_q^{\otimes n}$ which extends $\Phi_1$. In~\cite[Proposition~6.7]{BR}, we showed that $\Phi_n$ affords isomorphisms
\begin{equation}\label{Alekseevmap}
\Phi_n\colon\ \Ll_{0,n} \ra \big(U_q^{\otimes n}\big)^{\rm lf},\qquad \Phi_n\colon\ \Mm_{0,n}\rightarrow \big(U_q^{{\otimes} n}\big)^{U_q},
\end{equation}
where $\big(U_q^{\otimes n}\big)^{\rm lf}$ is the set of ${\rm ad}_n^r$-locally finite elements of $U_q^{{\otimes} n}$. We call $\Phi_n$ the {\it Alekseev map}; we do not recall here the definition of $\Phi_n$, for we will not use it. It is a key argument of the proof of~\eqref{Alekseevmap} that the set of locally finite elements of $U_q^{{\otimes} n}$ for $({\rm ad}^r)^{\otimes n}\circ \Delta^{(n)}$ coincides with $\big(U_q^{\rm lf}\big)^{\otimes n}$; this follows from the main result of \cite{KLNY}. Using that the map
\begin{equation}\label{Phipsi}
\psi_n = \Phi_n \circ \big(\Phi_1^{-1}\big)^{\otimes n}\colon\ \big(U_q^{\rm lf}\big)^{\otimes n}\ra \big(U_q^{\otimes n}\big)^{\rm lf}
\end{equation}
intertwines the actions $({\rm ad}^r)^{\otimes n}\circ \Delta^{(n-1)}$ and ${\rm ad}_n^r$, we deduced that ${\rm Im}(\Phi_n) = \big(U_q^{\otimes n}\big)^{\rm lf}$.

\begin{Remark}We have \smash{$\big(U_q^{\rm lf}\big)^{\otimes n} \!\ne\! \big(U_q^{\otimes n}\big)^{\rm lf}$} and in fact there is not even an inclusion.~Indeed, let $\Omega=\big(q-q^{-1}\big)^2FE+qK+q^{-1}K^{-1}$ be the Casimir element of $U_q({\mathfrak{sl}_2})$. We trivially have~\smash{$\Delta(\Omega)\in \big(U_q^{\otimes 2}\big)^{\rm lf}$} but
\[\Delta(\Omega)=\big(q-q^{-1}\big)^2\big(K^{-1}E\otimes FK+F\otimes E\big)+
\Omega\otimes K+ K^{-1}\otimes \Omega-\big(q+q^{-1}\big)K^{-1}\otimes K\]
and therefore \smash{$\Delta(\Omega)\notin \big(U_q^{\rm lf}\big)^{\otimes 2}$}, since \smash{$K\notin U_q^{\rm lf}$} (see, e.g., Theorem \ref{JLteo1}\,(2)). This reflects the fact that \smash{$U_q^{\rm lf}$} is only a right coideal of $U_q$ (and not a subcoalgebra).
\end{Remark}
As in Remark \ref{remsocle}, denote by $C(\mu)$, $\mu\in P^+$, the linear subspace of $\Ll_{0,1}$ generated by the matrix coefficients of $V_\mu$. For every tuple $[\mu] = (\mu_1,\dots,\mu_n)\in P_+^n$ put
\begin{equation}\label{defCmu}
C([\mu]) =C(\mu_1)\otimes \dots \otimes C(\mu_n).
\end{equation}
Then \smash{$ {\mathcal L}_{0,n}=\bigoplus_{[\mu]\in P_+^n} C({[\mu]})$}. Each space $C({[\mu]})$ is a finite-dimensional $U_q$-module under the action ${\rm coad}^r_n$, whence it is completely reducible. Therefore, $
\Ll_{0,n} = \Mm_{0,n} \oplus I$ as $U_q$-modules, where $I$ is the sum of nontrivial isotypical components of $\Ll_{0,n}$. The $\mc(q)$-linear projection map
\begin{equation}\label{Reynoldsdef}
\Rr\colon\ {\mathcal L}_{0,n}\ra \Mm_{0,n}, \qquad {\rm Ker}(\Rr)=I
\end{equation}
is called the {\it Reynolds operator}. For all $\alpha\in \Mm_{0,n}$, $\beta \in \Ll_{0,n}$ it satisfies $
\Rr(\alpha \beta) = \alpha \Rr(\beta)$.
This property will be crucial in the sequel, so let us recall a (classical) proof of it. We can write~${\beta = \Rr(\beta) + \gamma}$ with $\gamma \in I$, and then we have to show $\alpha \gamma \in I$. We can reduce to the case where $\gamma$ is contained in a simple summand $V$ of $I$. Multiplication by the invariant element $\alpha$ yields a surjective map $V \ra \alpha V$, which is a morphism of $U_q$-modules. Since $V$ is simple, it is either the $0$ map, or an isomorphism. In either cases it follows $\alpha V \subset I$ (in fact the first case cannot happen, for $\Ll_{0,n}$ has no nontrivial zero divisors, as explained after \eqref{Zinv}).

We can formulate the Reynolds operator in the following way. Recall that $\Oo_q$ has a unique left (or right, or $2$-sided) Haar integral, that is a linear map $h\colon \Oo_q\ra \mc(q)$ such that
\[%\label{Haardef}
h(1)=1 \qquad \text{and}\qquad ({\rm id}\otimes h)\Delta(\alpha) = h(\alpha)1,\qquad \forall \alpha\in \Oo_q.
\]
(See, e.g., \cite[Proposition 13.3.6]{CP}.) It vanishes on all matrix coefficients except the one of the trivial representation, to which it gives the value $1$. Denote by $\Delta_\Ll\colon {\mathcal L}_{0,n} \ra {\mathcal L}_{0,n}\otimes \Oo_q$ the right coaction dual to the action ${\rm coad}_n^r$ of $U_q$ on ${\mathcal L}_{0,n}$. Then, in analogy with the formula of the averaging operator $\mathcal{C}^\infty(G) \ra \mathcal{C}^\infty(G)^G$, $f\ra [f] = \int_G f\big(g^{-1} \cdot g\big) {\rm d}\mu(g)$, for a locally compact group $G$ with Haar measure ${\rm d}\mu(g)$, it is straightforward that
\begin{equation}\label{formuleRh}
\Rr = ({\rm id}\otimes h)\Delta_\Ll .
\end{equation}
Note that the complete reducibility of $\Ll_{0,n}$ discussed after \eqref{defCmu} follows also from Theorem~\ref{JLteo1}\,(1), since by \eqref{Phipsi} we have an isomorphism of $U_q$-modules
\[%\label{phigplus}
{\mathcal L}_{0,n} \stackrel{\Phi_n}{\longrightarrow} \big(U_q({\mathfrak g})^{\otimes n}\big)^{\rm lf} \stackrel{\psi_n^{-1}}{\longrightarrow}U_q^{\rm lf}({\mathfrak g})^{\otimes n},
\]
where ${\rm lf}$ means respectively locally finite for the action ${\rm ad}_n^r$ of $U_q({\mathfrak g})$ on $U_q({\mathfrak g})^{\otimes n}$, and locally finite for the action ${\rm ad}^r$ of $U_q({\mathfrak g})$ on $U_q({\mathfrak g})$. An explicit basis of $\Mm_{0,n}$ is described in \cite[Pro\-po\-sition~6.22]{BR}.

Finally, let us point out here two important consequences of~\eqref{Alekseevmap}. First, $\Phi_n$ yields isomorphisms between centers, $\mathcal{Z}(\Ll_{0,n}) \cong \mathcal{Z}(U_q)^{\otimes n}$ and \smash{$\mathcal{Z}\big(\Ll_{0,n}^{U_q}\big) \cong \mathcal{Z}\big(\big(U_{q}^{\otimes n}\big)^{U_q}\big)$}, where one can show that \cite[Lemma 6.29]{BR}
\begin{equation}\label{Zinv}
\mathcal{Z}\big(\big(U_{q}^{\otimes n}\big)^{U_q}\big)\cong \Delta^{(n)}(\mathcal{Z}(U_q))\bigotimes_{{\mathbb C}(q)}\mathcal{Z}(U_q)^{\otimes n}.
\end{equation}
Second, $\Ll_{0,n}$ (and therefore $\Mm_{0,n}$) has no nontrivial zero divisors because of the isomorphisms \smash{$\Phi_n\colon \Ll_{0,n} \ra \big(U_q^{\otimes n}\big)^{\rm lf}\subset U_q^{\otimes n}$} and $U_q^{\otimes n}\cong U_q\big(\mathfrak{g}^{\oplus n}\big)$, and the fact that $U_q\big(\mathfrak{g}^{\oplus n}\big)$ has no nontrivial zero divisors (proved, e.g., in \cite{DC-K}).

\subsection{Integral forms and specializations}\label{INTEG} Let $A=\mc\big[q,q^{-1}\big]$. We call integral form of a (Hopf) $\mc(q)$-algebra $H$ a (Hopf) $A$-subalgebra~${}_AH$ such that the canonical map ${}_AH\bigotimes_A \mc(q) \ra H$ is an isomorphism. Note that the standard notion of integral form of $\mc(q)$-algebra uses $\mz\big[q,q^{-1}\big]$ instead of $\mc\big[q,q^{-1}\big]$; our choice is made for simplicity ($\mc\big[q,q^{-1}\big]$ is a principal ideal domain, whereas $\mz\big[q,q^{-1}\big]$ is not).

\subsubsection{Definitions}\label{defintform} The {\it unrestricted} integral form of $U_q$ is the $A$-subalgebra $U_A = U_A(\mathfrak{g})$ introduced by De Concini--Kac--Procesi in \cite[Section 12]{DC-K-P2} (and in a differently normalized form in \cite{DC-K,DC-K-P1}). It is the smallest $A$-subalgebra of $U_q$ which contains the elements ($i=1,\dots,m$)
\begin{equation}\label{normgen0}
\bar{E}_i = \big(q_i-q_i^{-1}\big)E_i,\qquad \bar{F}_i = \big(q_i-q_i^{-1}\big)F_i,\qquad L_i,\qquad L_i^{-1}
\end{equation}
and is stable under the action of $\mathcal{B}(\mathfrak{g})$ given by the Lusztig automorphisms (see \eqref{rootvectdef}). Recall the root vectors $E_{\beta_k}$, $F_{\beta_k}$ defined in \eqref{rootvectdef}. Let us put $q_\beta := q^{(\beta,\beta)/2}$. The algebra $U_A$ is a free $A$-module with basis the monomials $\bar{E}_{\beta_1}^{p_1}\cdots \bar{E}_{\beta_N}^{p_N}K_\lambda \bar{F}_{\beta_N}^{n_N}\cdots \bar{F}_{\beta_1}^{n_1}$, where $\lambda \in P$ and we set
\[\bar{E}_{\beta_k}= \big(q_{\beta_k}-q_{\beta_k}^{-1}\big)E_{\beta_k},\qquad \bar{F}_{\beta_k}= \big(q_{\beta_k}-q_{\beta_k}^{-1}\big)F_{\beta_k}.\]
We denote \smash{$U_A^{\rm lf} := U_A \cap U_q^{\rm lf}$}.
The unrestricted integral form of \smash{$U_q^{\rm ad}$} is defined similarly, as the smallest $A$-subalgebra $U_A^{\rm ad} \subset U_A$ which contains the elements $\bar{E}_i$, $\bar{F}_i$ and $K_i^{\pm 1}$, for $i=1,\dots,m$, and is stable under the Lusztig action of $\mathcal{B}(\mathfrak{g})$.

For $\beta$ a positive root, we define the divided powers
\[E_{\beta}^{(k)} = \frac{E_\beta^k}{[k]_{q_\beta}!} , \qquad F_{\beta}^{(k)} = \frac{F_\beta^k}{[k]_{q_\beta}!}, \qquad k\in \mn.\]
The Lusztig {\it restricted} integral form of $U_q^{\rm ad}$ \cite{Lusztig2,Lusztig} (see also \cite[Chapter 9.3]{CP}) is the $A$-sub\-algebra~$U_A^{\rm res}$ generated by the elements ($i=1,\dots,m$, $k\in \mn^*$)
\[E_i^{(k)}= \frac{E_i^k}{[k]_{q_i}!},\qquad F_i^{(k)}= \frac{F_i^k}{[k]_{q_i}!} ,\qquad K_i ,\qquad K_i^{-1}.\]
The algebra $U_A^{\rm res}$ is a free $A$-module with Poincar\'e--Birkhoff--Witt (PBW) basis
\[%\label{PBWUAres}
E_{\beta_1}^{(p_1)}\cdots E_{\beta_N}^{(p_N)}\prod_{i=1}^m K_i^{\sigma_i}[ K_i; t_i]_{q_i} F_{\beta_N}^{(n_N)}\cdots F_{\beta_1}^{(n_1)},
\]
where $\sigma_i\in \{0,1\}$, $n_i,p_i, t_i\in \mn$, and we set $[ K_i; 0]_{q_i}:=1$ and
\[[ K_i; t ]_{q_i}= \prod_{s=1}^t \frac{K_iq_i^{-s+1}-K_i^{-1}q_i^{s-1}}{q_i^s-q_i^{-s}}.\]
The integral forms $U_A(\mathfrak{h})$, $U_A(\mathfrak{b}_\pm)$ and $U_A^{\rm res}(\mathfrak{h})$, $U_A^{\rm res}(\mathfrak{b}_\pm)$ associated to the subalgebras $\mathfrak{h}$,~$\mathfrak{b}_\pm \subset \mathfrak{g}$ are the subalgebras of $U_A$ and $U_A^{\rm res}$, respectively, defined in the obvious way. For instance, the~``Cartan'' subalgebra $U_A^{\rm res}(\mathfrak{h}) = U_q(\mathfrak{h})\cap U_A^{\rm res}$ is generated as a $A$-module by the elements $\prod_{i=1}^m K_i^{\sigma_i}[ K_i; t_i]_{q_i}$.

Denote by $\mathcal{C}_A$ the category of $U_A^{\rm res}$-modules of {\it type $1$}, i.e., free $A$-modules of finite rank which have a basis where the elements $K_i$ act diagonally with eigenvalues of the form $q_i^k$, $k\in \mz$ (in general, finiteness of the rank imposes eigenvalues of the form $\pm q_i^k$, $k\in \mz$). The category~$\mathcal{C}_A$ is a~rigid and tensor category. It is not semisimple, and this makes the study of $\mathcal{C}_A$ a~complicated task; for this, see \cite{BR}, and Section \ref{canbasemodqg} below. Every type $1$ finite-dimensional simple $U_q$-module $V_\mu$,~$\mu \in P_+$, has a $U_A^{\rm res}$-invariant full $A$-sublattice, that we denote by ${}_AV_\mu$. These~$U_A^{\rm res}$-modules form the simple objects of $\mathcal{C}_A$. Moreover, $\mathcal{C}_A \otimes \mc\big[q^{1/D},q^{-1/D}\big]$ is a ribbon category (see Section~\ref{intdual}).

The {\it integral} quantum function Hopf algebra $\Oo_A=\Oo_A(G)$ is the (type $1$) restricted dual of~$U_A^{\rm res}$, that is, the $A$-span of the matrix coefficients $x\mapsto v^i(\pi_V(x)v_i)$, $x\in U_A^{\rm res}$, for every module~$V$ in $\mathcal{C}_A$, where $(v_i)$ is an $A$-basis of $V$ and $\big(v^i\big)$ the dual $A$-basis of the dual module~$V^*$ (compare with the definition of $\Oo_q$). We can also regard $\Oo_A$ as the set of $A$-linear maps~${f\colon U_A^{\rm res}\ra A}$ such that ${\rm Ker}(f)$ contains a cofinite two sided ideal $I$, and $ \prod_{s=-r}^r (K_i - q_i^s) \in I$ for some $r\in \mathbb{N}$ and every $i$. Because of the inclusions of $U_A^{\rm res}(\mathfrak{h})$, $U_A^{\rm res}(\mathfrak{n}_\pm)$, $U_A^{\rm res}(\mathfrak{b}_\pm)$ in $U_A^{\rm res}$, there are Hopf epimorphisms from $\Oo_A$ to the $A$-duals of these subalgebras, that we denote by $\Oo_A(T_G)$, $\Oo_A(U_\pm)$ and $\Oo_A(B_\pm)$, respectively.

The algebra $\Oo_A$ has been introduced by Lusztig in \cite{Lusztig2,Lusztig}. It is an integral form of $\Oo_q$, so $\Oo_q = \Oo_A \bigotimes_A \mc(q)$.

$\Oo_A$ is also the restricted dual of the integral form $\Gamma=\Gamma(\mathfrak{g})$ of \smash{$U_q^{\rm ad}$} introduced by De Concini--Lyubashenko in \cite[Sections 2 and~3]{DC-L}; $\Gamma$ is the $A$-subalgebra of \smash{$U_q^{\rm ad}$} generated by the elements ($i=1,\dots,m$)
\[E_i^{(k)}= \frac{E_i^k}{[k]_{q_i}!} ,\qquad F_i^{(k)}= \frac{F_i^k}{[k]_{q_i}!} ,\qquad ( K_i; t )_{q_i}= \prod_{s=1}^t \frac{K_iq_i^{-s+1}-1}{q_i^s-1},\qquad K_i^{-1},\]
where $k\in \mn$, $t\in \mn$ (setting $( K_i; 0)_{q_i}=1$ by convention). Note that the definition of $\Gamma$ is less symmetric than that of $U_A^{\rm res}$. However, $\Gamma$ contains the elements $K_i$, and the commutation relations between the generators \smash{$E_i^{(k)}$}, \smash{$F_i^{(k)}$} imply that the symmetrized elements $[ K_i; t ]_{q_i}$ belong to $\Gamma$. In fact, let us denote $\Gamma(\mathfrak{h}) = U_q(\mathfrak{h}) \cap \Gamma$ and $\Gamma(\mathfrak{b}_\pm) = U_q(\mathfrak{b}_\pm) \cap \Gamma$. It is proved in \cite[Theorem 3.1]{DC-L} that $\Gamma(\mathfrak{h})$ contains $U_A^{\rm res}(\mathfrak{h})$ and that the elements \smash{$ \prod_{i=1}^m K_i^{-\sigma(t_i)}(K_i; t_i)_{q_i}$, $t_i\in \mn$}, where $\sigma(t)$ is the integer part of $t/2$, is an $A$-basis of $\Gamma(\mathfrak{h})$. A PBW basis of $\Gamma$ is formed by the monomials
\[%\label{PBWAbasisGamma}
E_{\beta_1}^{(p_1)}\cdots E_{\beta_N}^{(p_N)}\prod_{i=1}^m K_i^{-\sigma(t_i)}(K_i; t_i)_{q_i} F_{\beta_N}^{(n_N)}\cdots F_{\beta_1}^{(n_1)}.
\]

The inclusion $U_A^{\rm res}\subset \Gamma$ is strict, for the elements $(K_i; t)_{q_i}$, $t\ne 0$, do not belong to $U_A^{\rm res}$. However, the restriction functor $\mathcal{C}_\Gamma \ra\mathcal{C}_A$ is obviously an equivalence, where $\mathcal{C}_\Gamma$ is the category of $\Gamma$-modules of {\it type $1$}, i.e., free $A$-modules of finite rank which have a basis where the elements~$K_i$ act diagonally with eigenvalues of the form $q_i^k$, $k\in \mz$. Therefore, we can identify the two categories, and $\Oo_A$ with the (type $1$) restricted dual of $\Gamma$. We will thus consider the $U_A^{\rm res}$-modules~${}_AV_\mu$, $\mu\in P_+$, equally as $\Gamma$-modules. We will sometimes use $\Gamma$ instead of $U_A^{\rm res}$ in order to make direct the connection with results of De Concini--Lyubashenko about the integral pairings $\pi_A^\pm$ considered in Section~\ref{intdual}.

The {\it integral} form $\Ll_{0,1}^A$ of $\Ll_{0,1}$ is defined as the $U_A^{\rm res}$-module $\Oo_A$ endowed with the product of~$\Ll_{0,1}$. The {\it integral} form \smash{$\Ll_{0,n}^A$} of $\Ll_{0,n}$ is the braided tensor product of $n$ copies of \smash{$\Ll_{0,1}^A$}; in particular, \smash{$\Ll_{0,n}^A = \Oo_A^{\otimes n}$} as $U_A^{\rm res}$-modules. That the products of $\Ll_{0,1}$ and $\Ll_{0,n}$ are well defined over $A$ was shown in \cite[Proposition~6.9]{BR}.

The {\it integral} quantum moduli algebra is
\[\Mm_{0,n}^A := \big(\Ll_{0,n}^A\big)^{U_A^{\rm res}} = \big(\Ll_{0,n}^A\big)^{U_A}.\]

Finally, given $q=\e\in \mc^{\times}$ we define the {\it specializations} $U_\e$, $\Gamma_\e$, $\Oo_\e$, $\Ll_{0,n}^{\e}$ and $\Mm_{0,n}^{A,\e}$ as the~$\mc$-algebras obtained by tensoring $U_A$, $\Gamma$, $\Oo_A$, \smash{$\Ll_{0,n}^A$} and \smash{$\Mm_{0,n}^A$} respectively with $\mc_\e$, the $A$-module~$\mc$ where $q$ acts by multiplication by $\e$. Each one can also be defined as the quotient by the ideal generated by $q-\e$. We find convenient to use the notations
\begin{equation}\label{notambig0}
\big(U_A^{\otimes n}\big)^{U_A}_\e:= \big(U_A^{\otimes n}\big)^{U_A}\bigotimes_A \mc_\e,\qquad \big(U^{\otimes n}\big)^{\rm lf}_\e:= \big(U_A^{\otimes n}\big)^{\rm lf}\bigotimes_A \mc_\e .
\end{equation}
Let us stress here that when $\e$ is a root of unity, taking the locally finite part and taking the specialization at $\e$ are non commuting operations. Indeed, as shown by Theorem \ref{DCKteo1} below, $U_\e$~is finite over $\mathcal{Z}_0(U_\e)$ and therefore all its elements are locally finite for ${\rm ad}^r$; on another hand~\smash{${U_\e^{\rm lf} = U_A^{\rm lf}\bigotimes_A \mc_\e}$} does not contain the elements~$L_i$.

Similarly, taking invariants and taking the specialization at $\e$ are non commuting operations when $\e$ is a root of unity: indeed, it is easily checked that in this case \smash{$\big(U_A^{\otimes n}\big)^{U_A}_\e$} and \smash{$\big(U_\e^{\otimes n}\big)^{U_\e}$}, or \smash{$\Mm_{0,n}^{A,\e} = \Mm_{0,n}^A \bigotimes_A \mc_\e$} and \smash{$\big(\Ll_{0,n}^\e\big)^{U_\e}$}, are distinct spaces. When $\e$ is a root of unity, we will not consider the algebras $\Mm_{0,n}^{A,\e}$ in this paper.

Arguments similar to those mentioned at the end of Section \ref{ALGdef} imply that the algebras $\Ll_{0,n}^A$,~$\Mm_{0,n}^A$ and $\Ll_{0,n}^\ep$, \smash{$\Mm_{0,n}^{A,\ep}$}, $\ep\in \mc^\times$, have no nontrivial zero divisors (see \cite[Propositions~6.11 and~6.30]{BR}).

\subsubsection{Canonical bases and modified quantum groups} \label{canbasemodqg} Because the category $\mathcal{C}_A$ is not semisimple, it is not clear from the above definition of $\Oo_A$ whether or not it is a finitely generated algebra, if $\Mm_{0,n}^A$ is a direct summand of the $A$-module $\Ll_{0,n}^A$, or if the projection map \eqref{projOq} may be refined to a morphism between underlying $A$-modules.

Such properties, using the appropriate formalism developed by Kashiwara--Lusztig, indeed hold true, and will play a key role later. We state them precisely in Proposition \ref{genOA}, Theorem~\ref{teoLuzstigOAinv} and Proposition \ref{teoLuzstigOAscinde}. These results are consequences of the existence of an $A$-basis of $\Oo_A$ with favourable properties, which implies in particular that $\Oo_A$ is a free $A$-module. In order to introduce this $A$-basis it is necessary to consider a variant of $U_q^{\rm ad}$ introduced by Lusztig~\cite{Lusztig}, called {\it modified quantum group}, and use the Kashiwara--Lusztig theory of canonical bases \cite{Kashiwara1,Kashiwara2,Kashiwara3,Lusztig}. We are going to recall the background material step by step.

The Lusztig {\it modified quantum group} is the $\mc(q)$-algebra $\dot{\mathbf{U}}$ obtained by replacing $U_q^{\rm ad}(\mathfrak{h})$ with the direct sum of infinitely many one-dimensional algebras, generated by orthogonal idempotents~$1_\lambda$ indexed by the elements $\lambda$ of the weight lattice $P$ \cite[Chapter~23]{Lusztig}. Namely, as a vector space \smash{$ \dot{\mathbf{U}} = \bigoplus_{\lambda',\lambda''\in P} {}_{\lambda'}\dot{\mathbf{U}}{}_{\lambda''}$}, where
\[{}_{\lambda'}\dot{\mathbf{U}}{}_{\lambda''} = U_q^{\rm ad}\bigg/\bigg(\sum_{\alpha\in Q} \big(K_\alpha -q^{(\alpha,\lambda')}\big)U_q^{\rm ad}+ \sum_{\alpha\in Q} U_q^{\rm ad}\big(K_\alpha -q^{(\alpha,\lambda'')}\big)\bigg).\]
Denote by $\pi_{\lambda',\lambda''}\colon U_q^{\rm ad} \ra {}_{\lambda'}\dot{\mathbf{U}}{}_{\lambda''}$ the canonical projection. The product of $\dot{\mathbf{U}}$ is given by \smash{$\pi_{\lambda_1',\lambda_1''}(s)\pi_{\lambda_2',\lambda_2''}(t) = \pi_{\lambda_1',\lambda_2''}(st)$} if $\lambda_1''=\lambda_2'$ and zero otherwise. Set $1_\lambda := \pi_{\lambda,\lambda}(1)$. The algebra \smash{$\dot{\mathbf{U}}$} has not unit, but the family $(1_\lambda)_{\lambda\in P}$ can be regarded as a substitute of it. Denote by $\Delta$ the collection of maps
\[
\Delta_{\lambda'_1,\lambda'_2,\lambda''_1,\lambda_2''} \colon {}_{\lambda'_1+\lambda'_2} \dot{\mathbf{U}} {}_{\lambda''_1+\lambda''_2} \ra {}_{\lambda'_1} \dot{\mathbf{U}} {}_{\lambda''_1}\otimes {}_{\lambda'_2} \dot{\mathbf{U}} {}_{\lambda''_2}
\]
 such that
\begin{equation}\label{defcoprod}
\Delta_{\lambda'_1,\lambda'_2,\lambda''_1,\lambda''_2} \pi_{\lambda'_1+\lambda'_2,\lambda''_1+\lambda''_2} = (\pi_{\lambda'_1,\lambda''_1}\otimes \pi_{\lambda'_2,\lambda''_2})\Delta_{U_q^{\rm ad}},
\end{equation}
where \smash{$\Delta_{U_q^{\rm ad}}$} is the coproduct of \smash{$U_q^{\rm ad}$}. We can regard $\Delta$ as a (categorically completed) coproduct \smash{$\Delta \colon \dot{\mathbf{U}} \ra \dot{\mathbf{U}} {}^{\hat{\otimes}2}$}. There is a natural structure of $U_q^{\rm ad}$-bimodule on $\dot{\mathbf{U}}$, defined by
\begin{equation}\label{bimodstr}
t'\pi_{\lambda',\lambda''}(s)t'' = \pi_{\lambda'+\nu',\lambda''-\nu''}(t'st'')
\end{equation}
for all \smash{$s\in U_q^{\rm ad}$} and all elements \smash{$t',t''\in U_q^{\rm ad}$} of respective weights $\nu'$, $\nu''$. This structure affords a triangular decomposition of $\dot{\mathbf{U}}$: given bases $\{b^\pm\}$ of \smash{$U_q^{\rm ad}(\mathfrak{n}_\pm)$}, the set of elements $b^+1_\lambda b^-$ (or~$b^-1_\lambda b^+$, or $b^+b^-1_\lambda$), where $\lambda\in P$, is a basis of $\dot{\mathbf{U}}$.

Given any \smash{$U_q^{\rm ad}$}-module $X$ of type $1$, and any weight subspace $X^\lambda\subset X$ of weight $\lambda\in P$, one can define the action of an element $u 1_\lambda \in \dot{\mathbf{U}}$, \smash{$u\in U_q^{\rm ad}$}, on $X$ as the projection onto $X^\lambda$ followed by the action of $u$. This way, one can identify the category $\mathcal{C}$ with the one of finite-dimensional {\it unital} $\dot{\mathbf{U}}$-modules, where unital means that all elements $1_\lambda$ act as $0$ but a finite number of them, and $ \sum_{\lambda\in P} 1_\lambda$ acts as the identity. It is proved in \cite[Section~29.5.1]{Lusztig}, that
\[%\label{Oqalt}
\Oo_q = \left\lbrace f\colon \dot{\mathbf{U}} \ra \mc(q) \, \bigg\vert \, \begin{matrix} \text{$f$ is $\mc(q)$-linear and vanishes on some}\\ \text{two-sided ideal of finite codimension of $\dot{\mathbf{U}}$} \end{matrix}\right\rbrace.
\]

There is an analogous realization of $\Oo_A$, of the form (see \cite[Sections~23.2 and~29.5.2]{Lusztig}, and~\cite{Lusztig3})
\[%\label{OAalt}
\Oo_A = \left\lbrace f\colon \dot{\mathbf{U}}_A \ra A \, \bigg\vert \, \begin{matrix} \text{$f$ is $A$-linear and vanishes on some}\\ \text{two-sided ideal of finite corank of $\dot{\mathbf{U}}_A$} \end{matrix}\right\rbrace ,
\]
where $\dot{\mathbf{U}}_A$ is the $A$-subalgebra of $\dot{\mathbf{U}}$ generated by the elements \smash{$E_i^{(k)}1_{\lambda}$} and \smash{$F_i^{(k)}1_{\lambda}$}, for all $i\in \{1,\dots,m\}$, $k\in \mn$ and $\lambda \in P$. It is a $U_A^{\rm res}$-subbimodule of $\dot{\mathbf{U}}$, and the coproduct restricts to a map \smash{$\Delta \colon \dot{\mathbf{U}}_A \ra \dot{\mathbf{U}}_A {}^{\hat{\otimes}2}$}. The above identification of the category $\mathcal{C}$ with the one of finite-dimensional unital $\dot{\mathbf{U}}$-modules yields an identification of the category $\mathcal{C}_A$ of $U_A^{\rm res}$-modules of type $1$ with the category of \smash{$\dot{\mathbf{U}}_A$}-modules of finite rank.

The key advantage of this realization of $\Oo_A$ is that $\dot{\mathbf{U}}_A$ can be equipped with a canonical $A$-basis $\dot{\mathbf{B}}$. The construction of $\dot{\mathbf{B}}$ is described in \cite[Chapter~25]{Lusztig}. It relies on the Kashiwara--Lusztig {\it canonical basis} of $U_A^{\rm res}(\mathfrak{n}_-)$. This last basis, denoted by $\mathbf{B}^-$, is defined in \cite[Chapter~14]{Lusztig}, and \cite{Kashiwara1} (a review can be found in \cite[Chapter~14]{CP}). It enjoys the following nice properties. Denote by ${}^-\colon \mc(q) \ra \mc(q)$ the field involution such that $\overline{q} = q^{-1}$, and by ${}^-\colon U_q^{\rm ad} \ra U_q^{\rm ad}$ the homomorphism of $\mc$-algebras such that
\[\bar{E}_i = E_i,\qquad \bar{F}_i = F_i,\qquad \bar{K}_\lambda = K_{-\lambda},\qquad \overline{fx} = \bar{f}\bar{x}\]
for all $f\in \mc(q)$, $x\in U_q^{\rm ad}$ ($\bar{E}_i$ and $\bar{F}_i$ above, which will not appear elsewhere, should not be confused with the normalized elements in \eqref{normgen0}). The map ${}^-$ yields a $\mc$-algebra homomorphism~${{}^-\colon \dot{\mathbf{U}} \ra \dot{\mathbf{U}}}$. Then, we have
\begin{enumerate}\itemsep=0pt
\item[(1)] the elements of $\mathbf{B}^-$ are weight vectors under the adjoint action of $U_q^{\rm ad}(\mathfrak{h})$;
\item[(2)] for every $b\in \mathbf{B}^-$, $\bar{b} =b$;
\item[(3)] for every $b,b'\in \mathbf{B}^-$, $bb' = \sum_{b''\in \mathbf{B}^-} N^{bb'}_{b''} b''$ where $N^{bb'}_{b''}\in \mz\big[q,q^{-1}\big]$;
\item[(4)] for every $b,b'\in \mathbf{B}^-$, $\Delta(b) = \sum_{b',b''\in \mathbf{B}^-} C_{b'b''}^b b' \otimes b''$ where $C_{b'b''}^{b}\in \mz\big[q,q^{-1}\big]$;
\item[(5)] for every $\mu\in P^+$, denoting by $v_\mu$ the highest weight vector of the $U_A^{\rm res}$-module ${}_AV_\mu$, the elements $bv_\mu$ which are non-zero, where $b\in \mathbf{B}^-$, form an $A$-basis of ${}_AV_\mu$.
\end{enumerate}
When $\mathfrak{g}$ is simply laced, the coefficients \smash{$N^{bb'}_{b''}$} and \smash{$C_{b'b''}^b$} belong to $\mn\big[q,q^{-1}\big]$ \cite[Theorem~14.3.13]{Lusztig}. In the case of $\mathfrak{g}={\mathfrak{sl}_2}$, the elements of $\mathbf{B}^-$ are just the divided powers \smash{$F^{(k)}$}, $k\in \mn$. Formulas in terms of PBW basis elements are known also for $\mathfrak{g}=\mathfrak{sl}_3$ and $\mathfrak{sl}_4$, and an algorithm exists in the general case (see \cite{deGraaf} and the references therein).

Correspondingly to $\mathbf{B}^-$, the set $\mathbf{B}^+ = \omega(\mathbf{B}^-)$ is a basis of $U_A^{\rm res}(\mathfrak{n}_+)$, where $\omega\colon U_q^{\rm ad} \ra U_q^{\rm ad}$ is the ($\mc(q)$-linear) {\it Cartan automorphism}, defined by
\[\omega(E_i)=F_i , \qquad \omega(F_i) = E_i , \qquad \omega(K_i)=K_i^{-1}\]
for $i=1,\dots,m$. The triangular decomposition of $\dot{\mathbf{U}}$ implies that the elements $b^+1_\lambda b'{}^-$, where~${b^+ \in \mathbf{B}^+}$, $b'{}^- \in \mathbf{B}^-$ and $\lambda\in P$, form a basis of $\dot{\mathbf{U}}$. They form in fact an $A$-basis of~$\dot{\mathbf{U}}_A$, and its elements are fixed by the involution ${}^-\colon \dot{\mathbf{U}} {}\ra{} \dot{\mathbf{U}}$.

Lusztig has constructed another $A$-basis of $\dot{\mathbf{U}}_A$, denoted $\dot{\mathbf{B}}$, and called the {\it canonical basis} of~$\dot{\mathbf{U}}_A$. It satisfies numerous properties that we now review. Its elements are denoted by~${b\,\lozenge_\lambda\, b'}$, where $b, b' \in \mathbf{B}^-$ and $\lambda\in P$, and are related to the elements $b^+b'{}^-1_\lambda$, where $b^+ := \omega(b)$ and~${b'{}^-:=b'}$, by a specific trigonal change of basis with coefficients in~$A$. Although we always have $b^+1_\lambda$, $b'{}^-1_\lambda \in \dot{\mathbf{B}}$, to our knowledge explicit formulas of the elements of $\dot{\mathbf{B}}$ as linear combinations of elements $b^+1_\lambda b'{}^-$ or $b'{}^- 1_\lambda b^+$ are known only for $\mathfrak{g}={\mathfrak{sl}_2}$ or $\mathfrak{sl}_3$ (see \cite[Section~25.3]{Lusztig} and~\cite{Cui}). In the former case, identifying $P$ with $\mz$ and $Q$ with $2\mz$ the canonical basis $\dot{\mathbf{B}}$ is formed by the elements
\[%\label{basesl2}
E^{(k)}1_{-n}F^{(l)}\quad \text{and}\quad F^{(l)}1_nE^{(k)}, \qquad k,l,n\in \mn,\qquad n\geq k+l,
\]
where $E^{(k)}1_{-n}F^{(l)} = F^{(l)}1_nE^{(k)}$ for $n=k+l$.

We are going to review Lusztig's construction of $\dot{\mathbf{B}}$, its canonical partition $\dot{\mathbf{B}} = \bigcup_{\lambda\in P_+} \dot{\mathbf{B}} [\lambda]$, the dual basis \smash{$\dot{\mathbf{B}} {}^*$}, and Kashiwara's approach to \smash{$\dot{\mathbf{B}} {}^*$} \cite{Kashiwara2,Kashiwara3}. The latter is stated in Theorem~\ref{canKashi} below. At first we need to recall the notions of based module and balanced triple; for details on these notions we refer to \cite[Chapter~27]{Lusztig} and \cite{Kashiwara2} (see also \cite{Kashiwara4}, \cite[Sections~3.15 and~3.16]{VY}, or \cite[Chapter~14]{CP} for overviews).

Denote by $\mathcal{A}_0\subset \mc(q)$ the ring of rational functions regular at $q=0$. By applying the involution ${}^-$, put $\mathcal{A}_\infty=\overline{\mathcal{A}}_0$.
Since $\mathcal{A}_0$ is the localization of $\mc[q]$ at $q=0$, we may regard $\mathcal{A}_\infty$ as the localization of $\mc\big[q^{-1}\big]$ at~${q=\infty}$.

Let us recall briefly the definition of crystal basis (see \cite{Kashiwara1}). Denote by $U_q^{\rm ad}(\mathfrak{g})_i$ the subalgebra of $U_q^{\rm ad}(\mathfrak{g})$ generated by $E_i$, $F_i$ and $K_i^{\pm 1}$; thus $U_q^{\rm ad}(\mathfrak{g})_i$ is isomorphic to $U_{q_i}({\mathfrak{sl}_2})$. Let $M$ be a~$U_q^{\rm ad}$-module of type $1$. Denote \smash{$M^\zeta$} the subspace of $M$ of weight $\zeta\in P$. For every $i=1,\dots, m$, we can regard $M$ as a $U_q^{\rm ad}(\mathfrak{g})_i$-module, so $M\cong \bigoplus_j V_{\lambda_j}$ for some simple $U_q^{\rm ad}(\mathfrak{g})_i$-modules $V_{\lambda_j}$. These being generated by primitive weight vectors, the PBW basis of $U_q^{\rm ad}(\mathfrak{g})_i$ yields
\[M = \bigoplus_{\zeta\in P} \bigoplus_{0\leq n\leq (\check{\alpha}_i, \zeta)} F_i^{(n)}\big({\rm Ker}(E_i)\cap M^\zeta\big).\]
The {\it Kashiwara operators} $\tilde{e}_i$, $\tilde{f}_i$ are the endomorphisms of $M$ defined by, for every $v\in {\rm Ker}(E_i)\cap M^\zeta$ and $0\leq n\leq (\check{\alpha}_i, \zeta)$,
\[%\label{Kashiop}
\tilde{f}_i\big(F_i^{(n)}v\big) = F_i^{(n+1)}v,\qquad \tilde{e}_i\big(F_i^{(n)}v\big) = F_i^{(n-1)}v.
\]
A {\it crystal basis of $M$ at $q=0$} consists of a pair $(\Ll, \mathcal{B})$, where
\begin{itemize}\itemsep=0pt
\item $\Ll$ is a free $\mathcal{A}_0$-sublattice of $M$ such that the canonical map $\Ll \bigotimes_{\mathcal{A}_0} \mc(q) \ra M$ is an isomorphism;
\item $\mathcal{B}$ is a basis of the $\mc$-vector space $\Ll/q\Ll$;
\item $ \Ll = \bigoplus_{\zeta \in P} \Ll^\zeta$ and \smash{$\mathcal{B} = \coprod_{\zeta \in P} \big(\mathcal{B}\cap \Ll^\zeta/q\Ll^\zeta\big)$}, where $\Ll^\zeta = \Ll\cap M^\zeta$;
\item for every $i=1,\dots, m$ the Kashiwara operators $\tilde{e}_i$, $\tilde{f}_i$ preserve $\Ll$, and the induced maps on~$\Ll/q\Ll$ send $\mathcal{B}$ into $\mathcal{B}\cup \{0\}$, and satisfy $b' = \tilde{f}_i(b)$ if and only if $b=\tilde{e}_i(b')$ for every~${b,b'\in B}$.
\end{itemize}
Crystal bases at $q=\infty$ are defined similarly, by replacing $\mathcal{A}_0$ with $\mathcal{A}_\infty$ and $q$ with $q^{-1}$.


 A {\it based module} consists of a pair $(M,B)$ where $M$ is a~$U_q^{\rm ad}$-module of type $1$ endowed with a~$\mc(q)$-basis $B$ such that the following conditions hold:
\begin{itemize}\itemsep=0pt
\item[(i)] For every weight $\zeta\in P$, the set $B\cap M^{\zeta}$ is a basis of the weight subspace $M^{\zeta}\subset M$.
\item[(ii)] The $A$-module ${}_A M$ generated by $B$ is stable under $U_A^{\rm res}$.


We will denote by $\Ll_M$ the $\mathcal{A}_0$-submodule of $M$ generated by $B$, and by $\bar{\Ll}_M$ the $\mathcal{A}_\infty$-submodule of $M$ generated by $B$.


\item[(iii)] The $\mc$-linear involution ${}^-\colon M\ra M$ defined by $\overline{fb} = \overline{f}b$ for all $f\in \mc(q)$ and $b\in B$ is compatible with the action of $U_q^{\rm ad}$ in the sense that $\overline{xm} = \bar{x}\bar{m}$ for all $x\in U_q^{\rm ad}$, $m\in M$.
\item[(iv)] The $\mathcal{A}_\infty$-submodule $\bar{\Ll}_M$ of $M$ together with the image of $B$ in $\bar{\Ll}_M/q^{-1}\bar{\Ll}_M$ forms a crystal basis of $M$ at $q=\infty$.
\end{itemize}
If $(M,B)$ is a based module, we will denote by $\overline{\mathcal{B}}$ the image of $B$ in $\bar{\Ll}_M/q^{-1}\bar{\Ll}_M$. From the notion of balanced triple that we recall now, denoting by $\mathcal{B}$ the image of $B$ in $\Ll_M/q\Ll_M$, we see that $(\Ll_M,\mathcal{B})$ is a crystal basis at $q=0$.

Indeed, consider more generally a $\mc(q)$-vector space $V$, finite-dimensional or not, a sub-$A$-module ${}_AV$, a sub-$\mathcal{A}_0$-module ${}_{\mathcal{A}_0}V$ and a sub-$\mathcal{A}_\infty$-module ${}_{\mathcal{A}_\infty}V$ satisfying the conditions (all isomorphisms being the canonical maps)
\[V\cong \mc(q) \bigotimes_A {}_AV ,\qquad V\cong \mc(q) \bigotimes_{\mathcal{A}_0} {}_{\mathcal{A}_0}V ,\qquad V\cong \mc(q) \bigotimes_{\mathcal{A}_\infty} {}_{\mathcal{A}_\infty}V.\]
Consider the $\mc$-vector space $E := {}_AV \cap {}_{\mathcal{A}_0}V \cap {}_{\mathcal{A}_\infty}V$. Then $({}_AV, {}_{\mathcal{A}_0}V,{}_{\mathcal{A}_\infty}V)$ is a {\it balanced triple} \cite{Kashiwara1,Kashiwara2} if the canonical maps
\begin{equation}\label{condbalanced}
A \bigotimes_\mc E \ra {}_AV,\qquad \mathcal{A}_0 \bigotimes_\mc E \ra {}_{\mathcal{A}_0}V ,\qquad \mathcal{A}_\infty \bigotimes_\mc E \ra {}_{\mathcal{A}_\infty}V
\end{equation}
are isomorphisms. Equivalently, $({}_AV, {}_{\mathcal{A}_0}V,{}_{\mathcal{A}_\infty}V)$ is balanced if and only if the canonical map $E \ra {}_{\mathcal{A}_0}V/q{}_{\mathcal{A}_0}V$ is an isomorphism, if and only if the canonical map $E \ra {}_{\mathcal{A}_\infty}V/q^{-1}{}_{\mathcal{A}_\infty}V$ is an isomorphism \cite[Lemma 2.1.1]{Kashiwara2}.

Given a based module $(M,B)$, the elements of $B$ are weight vectors and $\overline{b}=b$ for every~${b\in B}$. Also, if an element $m\in {}_A M$ satisfies $\overline{m}=m$ and $m\in B + q^{-1}\bar{\Ll}_M$, then $m\in B$ (see \cite[Section~27.1.5]{Lusztig} for details on this fact). It follows that the canonical quotient map
\begin{equation}\label{basedbalancediso} {}_A M \cap \Ll_M \cap \bar{\Ll}_M \ra \bar{\Ll}_M/q^{-1}\bar{\Ll}_M
\end{equation}
is an isomorphism of $\mc$-vector spaces. This provides another way of viewing based modules: by \eqref{basedbalancediso}, $\big({}_A M, \Ll_M,\bar{\Ll}_M\big)$ is a balanced triple, and by \eqref{condbalanced} the $A$-lattice ${}_A M$ is completely determined by the crystal base $\big(\bar{\Ll}_M,\overline{\mathcal{B}}\big)$. We will say that $\big(\bar{\Ll}_M,\overline{\mathcal{B}}\big)$ (or the corresponding crystal base $(\Ll_M,\mathcal{B})$ at $q=0$) is {\it melted into} the based module $(M,B)$.

We will indifferently apply the notion of based module to finite-dimensional unital $\dot{\mathbf{U}}$-modules, since they are equivalent to $U_q^{\rm ad}$-modules of type $1$.


Every module $V_\mu$, $\mu\in P^+$, supports a structure of based module (see \cite[Section~14.4.10]{Lusztig} and~\cite{Kashiwara1}); the corresponding basis, called {\it canonical basis} and that we will denote by $\underline{\mathbf{B}}_{\mu}$, is formed by the elements $bv_{\mu}\in {}_AV_\mu$ which are non-zero, where $b\in \mathbf{B}^-$ and $v_{\mu}$ is the canonical highest weight vector of $V_{\mu}$, corresponding to the coset of $1\in U_q^{\rm ad}(\mathfrak{n}_-)$ in the Verma module construction of $V_{\mu}$. Note that the involution $\ \bar{}\ \colon V_{\mu}\ra V_{\mu}$ defined by (iii) above is indeed an automorphism, for the space $V_\mu$ with action of $U_q^{\rm ad}$ defined by $x\cdot v := \bar{x}v$, for all $x\in U_q^{\rm ad}$, $v\in V_\mu$, has highest weight $\mu$, and is thus isomorphic to $V_\mu$ via the map $\ \bar{}\ $. The crystal base \smash{$\big(\Ll_\mu^{\rm low},\mathcal{B}_\mu^{\rm low}\big)$} at $q=0$ is formed by the $\mathcal{A}_0$-sublattice $\Ll_\mu^{\rm low}$ of $V_{\mu}$ generated by $\underline{\mathbf{B}}_{\mu}$ (which is eventually the same as the $\mathcal{A}_0$-sublattice generated by the vectors of the form \smash{$\tilde{f}_{i_1}\circ \dots \circ \tilde{f}_{i_k}(v_\mu)$}, where $i_1,\dots ,i_k \in \{1,\dots,m\}$), and $\mathcal{B}_\mu^{\rm low}$ is the set of non-zero images of these vectors in \smash{$\Ll_\mu^{\rm low}/q\Ll_\mu^{\rm low}$}.

There is the following uniqueness result \cite[Theorem 3]{Kashiwara1}.

\begin{teo}\label{Kashiuniqueteo} Let $M$ be a $U_q^{\rm ad}$-module of type $1$, and $(\Ll, \mathcal{B})$ a crystal base at $q=0$ of~$M$. Then there exists a $\mc(q)$-isomorphism $M\ra \bigoplus_j V_{\lambda_j}$ by which $(\Ll, \mathcal{B})$ is $\mathcal{A}_0$-isomorphic to \smash{$ \bigoplus_j \big(\Ll_{\lambda_j}^{\rm low}, \mathcal{B}_{\lambda_j}^{\rm low}\big)$}.
\end{teo}

The based modules form a category. Given based modules $(M,B)$ and $(M',B')$, a morphism of $U_q^{\rm ad}$-modules $f\colon M\ra M'$ is a morphism of based modules if{\samepage
\begin{itemize}\itemsep=0pt
\item[(a)] $f(b)\in B'\cup\{0\}$ for any $b\in B$;
\item[(b)] $B\cap {\rm Ker}(f)$ is a basis of ${\rm Ker}(f)$.
\end{itemize}}%

\pagebreak

\noindent
The direct sum of based modules $(M,B)$ and $(M',B')$ is a based module $(M \oplus M',B\cup B')$; and a submodule $M'$ of a~based module $(M,B)$ spanned over $\mc(q)$ by a subset $B'$ of $B$ forms a~based module $(M',B')$. The quotient module $M/M'$ together with the image of $B\setminus B'$ is then a based module.

The tensor product of based modules $(M,B)$, $(M',B')$ is also defined. Namely, consider the~$\mc$-linear map $\Psi\colon M\otimes M' \ra M\otimes M'$ defined by
\[
\Psi(m\otimes m') = \hat{R}^{-1}(\bar{m}\otimes \bar{m}'),
\] where $\hat{R} = \Theta^{-1}R$, see \eqref{Rmatfact} (note that, as we use the coproduct opposite to \cite{Lusztig} our quasi-$R$-matrix is $\hat{R}^{-1}$). It can be checked that $\Psi$ is an involution compatible with the action of $\dot{\mathbf{U}}$ in the sense of (iii) above in the definition of based module. Moreover, denote by $\Ll_{M,M'}$ the~$\mc[q^{-1}]$-submodule of $M\otimes M'$ spanned by the basis elements~$b\otimes b'$, where $b\in B$, $b'\in B'$. It is shown in \cite[Section~27.3]{Lusztig}, that for every pair $(b,b')\in B\times B'$ there is a unique element $b\,\lozenge\, b'\in \Ll_{M,M'}$ such that
\begin{itemize}\itemsep=0pt
\item[(a)] $\Psi(b\,\lozenge\, b')=b\,\lozenge\, b'$,
\item[(b)] $b\,\lozenge\, b' - b\otimes b'\in q^{-1} \mathcal{\Ll}_{M,M'}$.
\end{itemize}
Moreover, $B_{\lozenge} = \{b\,\lozenge\, b', b\in B, b'\in B'\}$ is a basis of $M\otimes M'$, a $\mc[q^{-1}]$-basis of $\Ll_{M,M'}$, a $\mc\big[q,q^{-1}\big]$-basis of the $\mc\big[q,q^{-1}\big]$-module ${}_A\Ll_{M,M'}$ of $M\otimes M'$ generated by the elements $b\otimes b'$, where $b\in B$, $b'\in B'$, and $(M\otimes M',B_{\lozenge})$ is a based module.

This construction of $B_{\lozenge}$ is associative. Since $(V_\mu, \underline{\mathbf{B}}_{\mu})$ is for every $\mu\in P_+$ a based module, it follows that any tensor product $M$ of a finite number of the simple modules $V_\mu$ is naturally a based module. The corresponding basis is called {\it the canonical basis} of $M$. These canonical basis have been computed explicitly in \cite{FK} in the case $\mathfrak{g}={\mathfrak{sl}_2}$.

Consider now the $U_q^{\rm ad}$-module ${}^{\omega}V_{\mu}$ with underlying space $V_{\mu}$, $\mu\in P_+$, and action defined by $x._{\omega} v := \omega(x)v$, for every \smash{$x\in U_q^{\rm ad}$} and $v\in V_{\mu}$ (as usual $\omega\colon U_q^{\rm ad} \ra U_q^{\rm ad}$ is the Cartan automorphism). Note that there are isomorphisms \smash{${}^{\omega}V_{\mu}\cong V_{-w_0(\mu)}\cong V_{\mu}^*$} (endowed with the standard left action of $U_q^{\rm ad}$). Let us denote by ${}^\omega v_\mu$ the vector $v_\mu$ regarded in ${}^{\omega}V_{\mu}$ (i.e., its canonical lowest weight vector), and by ${}^\omega\underline{\mathbf{B}}_{\mu}:=\{b._{\omega}{}^\omega v_{\mu},\, b\in \mathbf{B}^+\}\setminus \{0\}$ its canonical basis; note that ${}^\omega\underline{\mathbf{B}}_{\mu} = \{\omega(b){}v_{\mu},\, b\in \omega(\mathbf{B}^-)\}\setminus \{0\} = \{b v_{\mu},\, b\in \mathbf{B}^-\}\setminus \{0\} = \underline{\mathbf{B}}_{\mu}$. Then ${}^{\omega}V_{\mu'}\otimes V_{\mu''}$ has the canonical basis $\underline{\mathbf{B}}_{\mu',\mu''} := \{\underline{b}'\,\lozenge\, \underline{b}'', \, \underline{b}'\in {}^\omega\underline{\mathbf{B}}_{\mu'},\, \underline{b}''\in \underline{\mathbf{B}}_{\mu''}\}$. Since $\underline{b}'\,\lozenge\, \underline{b}''$ is canonically determined by the elements $b' ,b''\in \mathbf{B}^-$ such that $\underline{b}' = \omega(b')._{\omega} {}^\omega v_{\mu'}$, $\underline{b}'' = b''v_{\mu''}$, following Lusztig we denote it by~${(b' \,\lozenge\, b'')_{\mu',\mu''}}$.

 Denote by $v_{w_0(\mu)}$ the canonical lowest weight vector of $V_{\mu}$, and by ${}^\omega v_{w_0(\mu)}$ the vector $v_{w_0(\mu)}$ regarded in ${}^{\omega}V_{\mu}$. It is a crucial observation that ${}^\omega v_{w_0(\mu')}\otimes v_{w_0(\mu'')}$ is a cyclic vector of ${}^{\omega}V_{\mu'}\otimes V_{\mu''}$ (see, e.g., \cite[Proposition~23.3.6]{Lusztig}; note that \smash{${}^\omega v_{w_0(\mu')}\otimes v_{w_0(\mu'')}$} plays the role of $\xi_{-\mu'}\otimes \eta_{\mu''}:= {}^\omega v_{\mu'}\otimes v_{\mu''}$ in \cite{Lusztig}, because we use opposite coproducts on $U_q^{\rm ad}$ but the factors ${}^{\omega}V_{\mu'}$ and $V_{\mu''}$ are ordered in the same way).

 We can now state the definition of the canonical basis $\dot{\mathbf{B}}$ of $\dot{\mathbf{U}}$: each element $u$ of $\dot{\mathbf{B}}$ belongs to $\dot{\mathbf{U}}_A 1_\zeta$ for some (unique) $\zeta\in P$, and it is then uniquely determined by the property that, for every $\mu'$, $\mu''\in P^+$ such that $w_0(\mu''-\mu')=\zeta$, we have
\begin{equation}\label{udef}
u({}^\omega v_{w_0(\mu')}\otimes v_{w_0(\mu'')}) = (b' \,\lozenge\, b'')_{\mu',\mu''}\end{equation}
for some \smash{$(b' \,\lozenge\, b'')_{\mu',\mu''} \in \underline{\mathbf{B}}_{\mu',\mu''}$} \cite[Section~25.2]{Lusztig}. We write \smash{$u = b' \, \lozenge_\zeta\, b''$}, and as in \cite{Lusztig3} we denote by \smash{$\dot{\mathbf{B}}_{\mu',\mu''}$} the finite subset of \smash{$\dot{\mathbf{B}}$} which is in bijection with $\underline{\mathbf{B}}_{\mu',\mu''}$ under the map $u\mapsto u({}^\omega v_{w_0(\mu')}\otimes v_{w_0(\mu'')})$. So
\begin{equation}\label{Bpointunion}
\dot{\mathbf{B}} =\bigcup_{\mu',\mu''\in P_+} \dot{\mathbf{B}}_{\mu',\mu''}.
\end{equation}
Note in particular that $\dot{\mathbf{B}}$ is formed by weight vectors for the left and right action of $U_q^{\rm ad}(\mathfrak{h})$ (defined as usual by~\eqref{bimodstr}).

In a sense, one can view $\dot{\mathbf{U}}$ as the projective limit of an inverse system formed by the \smash{$\big(U_q^{\rm ad}\otimes U_q^{\rm ad}\big)$}-modules \smash{${}^{\omega}V_{\mu'}\otimes V_{\mu''}$}, where $\mu',\mu''\in P^+$; then $\dot{\mathbf{B}}$ is the basis resulting from the corresponding inverse system of basis $\{\dot{\mathbf{B}}_{\mu',\mu''}\}_{\mu',\mu''}$.


Lusztig has produced a partition of $\dot{\mathbf{B}}$ as follows. First, consider the situation of a based module $(M,B)$. For every $\lambda\in P_+$ denote by $M[\lambda]$ the sum of the simple submodules of $M$ isomorphic to $V_\lambda$ (i.e., its isotypical component). Set
\begin{equation}\label{decbaseisospace}
M[\geq\lambda] = \bigoplus_{\lambda'\geq \lambda} M[\lambda'].
\end{equation}
Then, for every base element $b\in B$ there is a unique $\lambda\in P_+$ such that $b\in M[\geq\lambda]$ and $\lambda$ is maximal with this property \cite[Section~27.2]{Lusztig}. Denote by $B[\lambda]$ the set of all $b\in B$ that give rise to $\lambda\in P_+$ in this way. Clearly, the sets $B[\lambda]$, $\lambda\in P_+$, form a partition of $B$.

Now, given $b\in \dot{\mathbf{B}}$, let $\zeta\in P$ be the unique weight such that $b \in \dot{\mathbf{U}}_A 1_\zeta$, and let ${\mu', \mu''\in P^+}$ be such that $w_0(\mu''-\mu')=\zeta$, and $(\check{\alpha}_i,\mu')$ is large enough for all $i=1,\dots,m$ so that $u({}^\omega v_{w_0(\mu')}\otimes v_{w_0(\mu'')})$ is non-zero. This element belongs to the canonical basis $\underline{\mathbf{B}}_{\mu',\mu''}$ of ${}^{\omega}V_{\mu'}\otimes V_{\mu''}$, and therefore to one of the subsets $\underline{\mathbf{B}}_{\mu',\mu''}[\lambda]$, for a unique $\lambda\in P_+$. It is a result that $\lambda$ does not depend on the choice of $\mu'$, $\mu''$ (see \cite[Section~29.1.1]{Lusztig}). Hence there is a well-defined map~${\dot{\mathbf{B}}\ra P_+}$,~$b\mapsto \lambda$. Denoting by $\dot{\mathbf{B}} [\lambda]$ the fiber of this map, we thus obtain a partition
\begin{equation}\label{partition} \dot{\mathbf{B}} = \coprod_{\lambda\in P_+}\dot{\mathbf{B}} [\lambda]. \end{equation}
The sets $\dot{\mathbf{B}} [\lambda]$ are called {\it $2$-sided cells}. They are finite sets and have the following remarkable properties. For every $\lambda\in P_+$ denote by~${\dot{\mathbf{U}} [\geq \lambda]}$ and $\dot{\mathbf{U}} [> \lambda]$ the subspaces of $\dot{\mathbf{U}}$ spanned by~\smash{$ \coprod_{\lambda'\geq \lambda} \dot{\mathbf{B}} [\lambda']$} and \smash{$ \coprod_{\lambda'> \lambda} \dot{\mathbf{B}}[\lambda']$} respectively. Then $\dot{\mathbf{U}} [\geq \lambda]$ (respectively~${\dot{\mathbf{U}} [> \lambda]}$) consists of the elements $u\in \dot{\mathbf{U}}$ such that if $u$ acts on \smash{$V_\mu$} by a non-zero linear map, then $\mu\geq \lambda$ (respectively~${\mu> \lambda}$) \cite[Lemmas~29.1.3 and~29.1.4]{Lusztig}. Both $\dot{\mathbf{U}} [\geq \lambda]$ and $\dot{\mathbf{U}} [> \lambda]$ are two-sided ideals of~$\dot{\mathbf{U}}$. Moreover, the algebra homomorphism $\pi_\lambda\colon \dot{\mathbf{U}} [\geq\lambda] \ra {\rm End}(V_\lambda)$ given by the $\dot{\mathbf{U}}$-module structure on $V_\lambda$ descends to an algebra and $U_q^{\rm ad}$-bimodule isomorphism (keeping the same notation) \cite[Proposition~29.2.2]{Lusztig}
\begin{equation}\label{pimuiso} \bar{\pi}_\lambda\colon\ \dot{\mathbf{U}} [\geq \lambda]/ \dot{\mathbf{U}} [> \lambda] \ra {\rm End}(V_\lambda) . \end{equation}

For instance, when $\mathfrak{g}={\mathfrak{sl}_2}$ the $2$-sided cell $\dot{\mathbf{B}} [n]$ associated to the simple $U_q^{\rm ad}({\mathfrak{sl}_2})$-module of type $1$ and dimension $n+1$ is the set of cardinality $(n+1)^2$ given by \cite[Section~29.4.3]{Lusztig}
\begin{equation}\label{2sidedcellsl2}
\dot{\mathbf{B}} [n] =\big\{E^{(k)}1_{-n}F^{(l)},\,n\geq k+l \big\} \cup \big\{F^{(l)}1_nE^{(k)},\,n\geq k+l\big\},
\end{equation}
with the identification $E^{(k)}1_{-n}F^{(l)} = F^{(l)}1_nE^{(k)}$ when $n=k+l$. As we are mainly interested in~$\Oo_A$, we are going to describe the dual partition of $\dot{\mathbf{B}} {}^*$, see Theorem~\ref{canKashi}. The duality with~\eqref{partition} is discussed after that theorem.

First, we follow the approach of Kashiwara \cite{Kashiwara2,Kashiwara3}. For every $\lambda\in P_+$, denote by $V_\lambda^r$ the dual space of $V_\lambda$ endowed with its natural structure of right $U_q^{\rm ad}$-module, defined by ${(f x)(v) = f(xv)}$
for every $f\in V_\lambda^r$, \smash{$x\in U_q^{\rm ad}$}, $v\in V_\lambda$.
Clearly, $V_\lambda^r$ is a simple module of highest weight $\lambda$.
 Let~\smash{$\varphi\colon U_q^{\rm ad} \ra U_q^{\rm ad}$} be the anti-automorphism of $\mc(q)$-algebra given by $\varphi(E_i) = F_i$, $\varphi(F_i) = E_i$, $\varphi(K_\lambda) = K_{\lambda}$. By using $\varphi$, any right $U_q^{\rm ad}$-module can be considered as a left $U_q^{\rm ad}$-module. In particular, by the Verma module construction of $V_\lambda$ it follows
\[V_\lambda^r \cong U_q^{\rm ad}\bigg/\bigg(\sum_{\mu\in P_+} \big(K_\mu -q^{(\lambda,\mu)}\big)U_q^{\rm ad}+ \sum_{i=1}^m E_i^{1+(\check{\alpha}_i,\lambda)}U_q^{\rm ad}\bigg),\]
and $\varphi$ affords an isomorphism of the right module $V_\lambda^r$ with the {\it left} module $V_\lambda$. We will denote by $f_\lambda$ the unique highest weight vector of $V_\lambda^r$ satisfying $\langle f_\lambda,v_\lambda\rangle = 1$.

The space $V_\lambda^r \otimes V_\lambda$ can be identified with ${\rm End}(V_\lambda)^*$, and thus acquires by duality a natural structure of $U_q^{\rm ad}$-bimodule \big(or equivalently left $U_q^{\rm ad} \otimes \big(U_q^{\rm ad}\big)^{\rm op}$-module\big); the left and right actions are given by
\begin{equation}\label{bimodaction}x(f \otimes v)y = fy \otimes xv
\end{equation}
for every $x,y\in U_q^{\rm ad}$, $f\in V_\lambda^r$, $v\in V_\lambda$. The space $V_\lambda^r \otimes V_\lambda$ also acquires by duality a natural ``upper" crystal structure over $U_q^{\rm ad} \otimes \big(U_q^{\rm ad}\big)^{\rm op}$, as we explain now. Denote by $\langle \ ,\ \rangle_\lambda \colon V_\lambda \times V_\lambda \ra \mc(q)$ the unique symmetric bilinear form such that
\begin{equation}\label{proppairinglambda}
\langle v_\lambda,v_\lambda\rangle_\lambda =1\qquad{\rm and}\qquad \langle \varphi(x)u,v\rangle_\lambda = \langle u, xv\rangle_\lambda
\end{equation}
for every $u,v\in V_\lambda$ and $x\in U_q^{\rm ad}$. Recall the crystal base $\big(\Ll_\mu^{\rm low},\mathcal{B}_\mu^{\rm low}\big)$ at $q=0$ introduced before Theorem \ref{Kashiuniqueteo}. In Kashiwara's terminology \cite{Kashiwara1, Kashiwara2}, the pair $\big(\Ll_{\lambda}^{\rm low},\mathcal{B}_\lambda^{\rm low}\big)$ is the {\it lower crystal base} of $V_\lambda$ at $q=0$. Applying the involution ${}^-\colon V_\lambda\ra V_\lambda$, one obtains the lower crystal base \smash{$\big(\overline{\Ll_{\lambda}^{\rm low}},\overline{\mathcal{B}_\lambda^{\rm low}}\big)$} at $q=\infty$. Because the canonical bases are determined by the crystal bases (see the discussion about \eqref{basedbalancediso}), we call $(V_\lambda,\underline{\mathbf{B}}_\lambda)$ the {\it lower} based module of $V_\lambda$, and $\underline{\mathbf{B}}_\lambda$ the {\it lower canonical basis} of $V_\lambda$.

Put
\begin{align}
& {}_AV_\lambda^{\rm up} := \{v\in V_\lambda, \langle v,{}_AV_\lambda\rangle_\lambda\subset A\},\qquad
\Ll_{\lambda}^{\rm up} := \big\{v\in V_\lambda, \langle v,\Ll_{\lambda}^{\rm low}\rangle_\lambda\subset \mathcal{A}_0\big\},\nonumber\\
&\overline{\Ll_{\lambda}^{\rm up}} := \big\{v\in V_\lambda, \langle v,\overline{\Ll_{\lambda}^{\rm low}}\rangle_\lambda\subset \mathcal{A}_\infty\big\}.\label{upmodule}
\end{align}
Then $\big({}_A\!V_\lambda^{\rm up},\Ll_{\lambda}^{\rm up},\overline{\Ll_{\lambda}^{\rm up}}\big)$ is a balanced triple \cite[Lemma 4.2.1]{Kashiwara2}. Denote by $\mathcal{B}_\lambda^{\rm up}\!$ the basis of~${\Ll_{\lambda}^{\rm up}\!/\!q\Ll_{\lambda}^{\rm up}}$ dual to $\mathcal{B}_\lambda^{\rm low}$ by the induced pairing $\langle\ ,\ \rangle_\lambda\colon \Ll_{\lambda}^{\rm up}/q\Ll_{\lambda}^{\rm up} \times \Ll_{\lambda}^{\rm low}/q\Ll_{\lambda}^{\rm low} \ra \mc$. The pair $\big(\Ll_{\lambda}^{\rm up},\mathcal{B}_\lambda^{\rm up}\big)$ is the {\it upper crystal base} of $V_\lambda$ at $q=0$. The weight spaces of the $\mathcal{A}_0$-modules $\Ll_{\lambda}^{\rm low}$ and $\Ll_{\lambda}^{\rm up}$ are related by
\begin{equation}\label{renormweightA} \big(\Ll_{\lambda}^{\rm up}\big)^\mu = q^{\frac{(\lambda,\lambda)}{2} - \frac{(\mu,\mu)}{2}} \big(\Ll_{\lambda}^{\rm low}\big)^\mu , \qquad \mu\in P. \end{equation}
Correspondingly, denoting $\big(\mathcal{B}_\lambda^{\rm up}\big)^\mu := \mathcal{B}_\lambda^{\rm up} \cap \big(\Ll_\lambda^{\rm up}\big)^\mu$ and $\big(\mathcal{B}_\lambda^{\rm low}\big)^\mu := \mathcal{B}_\lambda^{\rm low} \cap \big(\Ll_\lambda^{\rm low}\big)^\mu$, we have (see~\cite{Kashiwara1} and \cite[equation~(4.2.9)]{Kashiwara2})
\[\big(\mathcal{B}_\lambda^{\rm up}\big)^\mu = q^{\frac{(\lambda,\lambda)}{2} - \frac{(\mu,\mu)}{2}} \big(\mathcal{B}_\lambda^{\rm low}\big)^\mu.\]
The $A$-module ${}_AV_\lambda^{\rm up}$ is characterized by the following two properties \cite[equations~(4.2.10)--(4.2.12)]{Kashiwara2}:
\begin{align*}
\big({}_AV_\lambda^{\rm up}\big)^\lambda & = \mc\big[q,q^{-1}\big]v_\lambda,\qquad
\big({}_AV_\lambda^{\rm up}\big)^\mu = \bigl\{v\in V_\lambda \mid U_A^{\rm res}\big(\mathfrak{n}^+\big)^{\lambda-\mu} v \in \mc\big[q,q^{-1}\big]v_\lambda\bigr\},
\end{align*}
where $U_A^{\rm res}\big(\mathfrak{n}^+\big)^{\gamma} = \big\{u\in U_A^{\rm res}\big(\mathfrak{n}^+\big) \mid \forall \nu\in P, \, K_\nu uK_\nu^{-1} = q^{(\nu,\gamma)}u\big\}$. Denote by $\underline{\mathbf{B}}_\lambda^{\rm up}$ the inverse image of $\mathcal{B}_\lambda^{\rm up}$ by the isomorphism \smash{${}_AV_\lambda^{\rm up}\cap \Ll_{\lambda}^{\rm up}\cap \overline{\Ll_{\lambda}^{\rm up}} \ra \Ll_{\lambda}^{\rm up}/q\Ll_{\lambda}^{\rm up}$}. By \eqref{condbalanced}, the set $\underline{\mathbf{B}}_\lambda^{\rm up}$ is a~basis of ${}_AV_\lambda^{\rm up}$; we call it the {\it upper canonical basis} of $V_\lambda$. In the appendix, we describe in details the ${\mathfrak{sl}_2}$ case.

Similarly, the right module $V_\lambda^r$ with its canonical basis $\underline{\mathbf{B}}_\lambda^r{}=\{f_\lambda b,\, b\in \mathbf{B}^+\}\setminus \{0\}$ has the lower crystal base $\big(\Ll_{\lambda}^r{}^{\rm low},\mathcal{B}_\lambda^r{}^{\rm low}\big)$, and it supports a balanced triple $\big({}_AV_\lambda^r{}^{\rm up},\Ll_{\lambda} ^r{}^{\rm up},\overline{\Ll_{\lambda} ^r{}^{\rm up}}\big)$ defined again by duality. We denote by $(\Ll_{\lambda}^r{}^{\rm up},\mathcal{B}_\lambda^r{}^{\rm up})$ and $\underline{\mathbf{B}}_\lambda^r{}^{\rm up}$ the corresponding crystal base and upper canonical basis of $V_\lambda^r$, respectively.

It follows that $\big({}_AV_\lambda^r{}^{\rm up}\bigotimes_A {}_AV_\lambda^{\rm up},\Ll_{\lambda}^r{}^{\rm up} \bigotimes_{\mathcal{A}_0}\Ll_{\lambda}^{\rm up},\overline{\Ll_{\lambda}^r{}^{\rm up}}\bigotimes_{\mathcal{A}_\infty}\overline{\Ll_{\lambda}^{\rm up}}\big)$ is a balanced triple; equivalently~${V_\lambda^r \otimes V_\lambda}$ with the bimodule structure \eqref{bimodaction} and the basis $\underline{\mathbf{B}}_\lambda^r{}^{\rm up}\otimes \underline{\mathbf{B}}_\lambda^{\rm up}$ is a based $\big(U_q^{\rm ad} \otimes (U_q^{\rm ad})^{\rm op}\big)$-module.

Denote again by $\langle\cdot,\cdot \rangle \colon \Oq \times \dot{\mathbf{U}} \ra \mc(q)$ the pairing of $U_q^{\rm ad}$-bimodules induced by the canonical pairing $\langle\ ,\ \rangle \colon \Oq \times U_q^{\rm ad} \ra \mc(q)$, and let $\Phi_\lambda\colon V_\lambda^r \otimes V_\lambda \ra \Oo_q$, $\lambda\in P_+$, be the ``matrix coefficient" map, i.e.,
\begin{equation}\label{pairingOAKashi}
\langle \Phi_\lambda (f \otimes v),x\rangle = \langle f, xv\rangle_\lambda
\end{equation}
for every $f\in V_\lambda^r$, $x\in U_q^{\rm ad}$, $v\in V_\lambda$. The map $\Phi := \bigoplus_{\lambda\in P_+} \Phi_\lambda$ is an isomorphism of $U_q^{\rm ad}$-bimodules, so let us use it to identify $\Oo_q$ with $\bigoplus_{\lambda\in P_+} V_\lambda^r \otimes V_\lambda$ (which is the content of the Peter--Weyl decomposition \eqref{decompdirectOq}). Define
\begin{alignat*}{3}
& \Ll(\Oo_q) = \bigoplus_{\lambda\in P_+} \biggl(\Ll_{\lambda}^r {}^{\rm up} \bigotimes_{\mathcal{A}_0} \Ll_{\lambda}^{\rm up}\biggr),\qquad && \mathcal{B}(\Oo_q) := \coprod_{\lambda\in P_+} \mathcal{B}_\lambda^r {}^{\rm up} \otimes \mathcal{B}_\lambda^{\rm up},& \\
& \overline{\Ll}(\Oo_q) = \bigoplus_{\lambda\in P_+} \biggl(\overline{\Ll_{\lambda}^r {}^{\rm up}} \bigotimes_{\mathcal{A}_\infty} \overline{\Ll_{\lambda}^{\rm up}}\biggr),\qquad && \overline{\mathcal{B}}(\Oo_q) := \coprod_{\lambda\in P_+} \overline{\mathcal{B}_\lambda^r {}^{\rm up}} \otimes \overline{\mathcal{B}_\lambda^{\rm up}}.&
\end{alignat*}
\begin{teo}\label{canKashi} \quad
\begin{itemize}\itemsep=0pt
\item[$(i)$] The triple $\big(\Oo_A,\Ll(\Oo_q),\overline{\Ll}(\Oo_q)\big)$ is balanced. Therefore, denoting by $G$ the inverse of the canonical map $\Oo_A \cap \Ll(\Oo_q) \cap \overline{\Ll}(\Oo_q)\ra \Ll(\Oo_q)/q\Ll(\Oo_q)$, we have
\[\Oo_A = \bigoplus_{b\in \mathcal{B}(\Oo_q)} A G(b).\]
\item[$(ii)$] The basis $G(\mathcal{B}(\Oo_q)):=\{G(b), \, b\in \mathcal{B}(\Oo_q)\}$ coincides with the dual canonical basis $\dot{\mathbf{B}} {}^*$, i.e., the elements $a^* \in \Oo_A$, for every $a\in \dot{\mathbf{B}}$, defined by $a^* (a') = \delta_{a,a'}$ for every $a'\in \dot{\mathbf{B}}$. Therefore, \[\Oo_A = \bigoplus_{b\in \dot{\mathbf{B}}} Ab^*.\]
\end{itemize}
\end{teo}
The statement (i) is \cite[Theorem~1]{Kashiwara2}, and (ii) is \cite[Theorem 10.1 and Proposition 10.2.2]{Kashiwara3} and \cite[Section~29.5]{Lusztig}. The basis $G(\mathcal{B}(\Oo_q))=\dot{\mathbf{B}} {}^*$ is called the {\it global basis}, or {\it canonical basis}, of~$\Oo_q$. The proof of Theorem \ref{canKashi}\,(ii) in \cite{Kashiwara3} (see also \cite{Kashiwara4}) exhibits an isomorphism of crystals over $U_q^{\rm ad} \otimes \big(U_q^{\rm ad}\big){}^{\rm op}$,
\begin{equation}\label{isocrystal}
\psi\colon\ \mathcal{B}(\Oo_q) \ra \mathcal{B}\bigl(\dot{\mathbf{U}}\bigr),
\end{equation}
where \smash{$\big(\Ll\bigl(\dot{\mathbf{U}}\bigr),\mathcal{B}\bigl(\dot{\mathbf{U}}\bigr)\big)$} is the crystal base of $\dot{\mathbf{U}}$ associated to the canonical basis $\dot{\mathbf{B}}$. The isomorphism $\psi$ satisfies $\langle G(b),G(b')\rangle =\delta_{\psi(b),b'}$ for every $b\in \mathcal{B}(\Oo_q)$, $b'\in \mathcal{B}\bigl(\dot{\mathbf{U}}\bigr)$. The unit $1$ of $\Oo_A$ is~$(1_0)^*$; the constant structures of $\Oo_A$ are studied in \cite{Lusztig,Lusztig3}.

{\bf The canonical basis of $\boldsymbol{\Oo_A}$ when $\boldsymbol{\mathfrak{g}={\mathfrak{sl}_2}}$}. Denote by $a$, $b$, $c$, $d$ the matrix coefficients in the canonical basis $(v_+,v_-:=Fv_+)$ of $V_1$, the simple $U_q^{\rm ad}({\mathfrak{sl}_2})$-module of type $1$ and dimension two, read from the top left to the bottom right. In that case of $V_1$ the upper canonical basis~$\underline{\mathbf{B}}_{1}^r{}^{\rm up}$ and \smash{$\underline{\mathbf{B}}_{1}^{\rm up}$} coincide with the lower ones (this is not true in general, see Example~\ref{V2V20}). The basis~$\dot{\mathbf{B}} {}^*({\mathfrak{sl}_2})$ is formed by the monomials $c^sa^pb^r$ where $p,r,s\in \mn$, and $c^sd^pb^r$ where $p,r,s\in \mn$ and $p>0$; this is stated in \cite[Proposition~9.1.1]{Kashiwara2} (in \cite[Proposition~1.3]{DC-L}, similar monomials are shown to form an $A$-basis of $\Oo_A({\rm SL}_2)$, but without reference to the canonical basis; see the comments before \eqref{PWrel} below). More precisely, recall the $2$-sided cells~\eqref{2sidedcellsl2}. We verified by a~tedious though straightforward computation that we have the duality pairing
\begin{gather*}
\begin{split}
& \big\langle c^sd^pb^r, E^{(i)}1_{-k}F^{(j)} \big\rangle = \delta_{p+r+s,k}\delta_{r,i}\delta_{s,j} ,\qquad \big\langle c^sd^pb^r, F^{(j)}1_{k}E^{(i)} \big\rangle =0, \\
& \big\langle c^sa^pb^r, E^{(i)}1_{-k}F^{(j)} \big\rangle = 0 ,\qquad \big\langle c^sa^pb^r, F^{(j)}1_{k}E^{(i)} \big\rangle =\delta_{p+r+s,k}\delta_{r,i}\delta_{s,j}.
\end{split}
\end{gather*}
Therefore, \begin{align*}\dot{\mathbf{B}} [n]^* :={}& \{ c^sa^pb^r, \, p,r,s\in \mn,\, p+r+s=n\} \\&\cup \{ c^{s}d^{p}b^{r},\, p,r,s\in \mn,\, p>0,\, p+r+s=n\}.\end{align*}
A description of $\dot{\mathbf{B}}^*$ in the case of $\mathfrak{g} = {\mathfrak{sl}_n}$ can be found in \cite{Du}. Moreover, denote by $V_n$ the simple~\smash{$U_q^{\rm ad}({\mathfrak{sl}_2})$}-module of type $1$ and dimension $n+1$, by $(v_k)$ the canonical basis of $V_n$, by~$\big(v^k\big)$ the dual basis, and by \smash{$\pi_n\colon \dot{\mathbf{U}}({\mathfrak{sl}_2}) \rightarrow {\rm End}(V_n)$} the representation morphism. By using the above pairing, it is readily checked that for every $0\leq l,m\leq n$, we have
\begin{gather}
v^l(\pi_n( \cdot) \ v_m)\nonumber\\
\qquad= \sum_{\substack{0\leq i,j,k\\ i+j\leq k\leq n\\j-i=l-m}} \hspace*{-0.5cm}\delta_{-k,n-2(m+j)} \left[\begin{matrix} m+j \\ j \end{matrix}\right]_{q} \left[\begin{matrix}n-m+i-j \\ i \end{matrix}\right]_{q}\big(E^{(i)}1_{-k}F^{(j)}\big)^*\nonumber\\
\phantom{\qquad=}{} + \sum_{\substack{0\leq i,j,k\\ i+j< k\leq n
\\j-i=l-m}}\delta_{k,n-2(m-i)} \left[\begin{matrix} m-i+j \\ j \end{matrix}\right]_{q} \left[\begin{matrix} n-m+i \\ i \end{matrix}\right]_{q}\big(F^{(j)}1_{+k}E^{(i)}\big)^*.\label{formulesl2}
\end{gather}
In particular, we see in this case of $\mathfrak{g}={\mathfrak{sl}_2}$ that in general the matrix coefficients of simple~\smash{$U_A^{\rm res}$}-modules of type $1$ are not elements of the dual canonical basis $\dot{\mathbf{B}} {}^*$. Moreover, these matrix coefficients do not form a basis of $\Oo_A$. For instance, it follows from \eqref{formulesl2} that the matrix of matrix coefficients of $V_2$ has the following form:
\begin{equation}\label{matrixcoefV2}
\begin{pmatrix}a^2 & [2]_qab & b^2 \\ ca & [2]_qbc +1 & db \\ c^2 & [2]_q cd & d^2 \end{pmatrix}.
\end{equation}
The matrix coefficient $v_0^* \otimes v_0$ being equal to $[2]_qbc +1$, this shows $bc$ cannot be expressed as a~linear combination over $A$ of matrix coefficients of simple modules.

{\bf The refined Peter--Weyl theorem.} Let us discuss the $U_A^{\rm res}$-bimodule structure of $\Oo_A$, and its relation with the partition \eqref{partition}. For every $\lambda\in P_+$, put
\begin{equation}\label{AClambda}
{}_A \overset{\raisebox{-1.5pt}{\scriptsize$\smallbullet$}}{C} (\lambda) := \bigoplus_{b\in \dot{\mathbf{B}}[\lambda]} Ab^*
\end{equation}
and
\[\Oo_A(\leq \lambda) := \bigoplus_{\lambda'\leq \lambda} {}_A \overset{\raisebox{-1.5pt}{\scriptsize$\smallbullet$}}{C} (\lambda') ,\qquad \Oo_A(< \lambda) := \bigoplus_{\lambda'< \lambda} {}_A \overset{\raisebox{-1.5pt}{\scriptsize$\smallbullet$}}{C} (\lambda').\]
In particular, in the ${\mathfrak{sl}_2}$ case the $A$-module ${}_A\! \overset{\raisebox{-1.5pt}{\scriptsize$\smallbullet$}}{C} (n\varpi_1)$ has basis $\dot{\mathbf{B}} [n]^*$ given above, of cardinality~${(n+1)^2}$.

Recall that $\dot{\mathbf{U}} [\geq \lambda]$ and $\dot{\mathbf{U}} [> \lambda]$ are two-sided ideals of $\dot{\mathbf{U}}$, and the algebra (whence $U_q^{\rm ad}$-bimodule) isomorphism \smash{$\bar{\pi}_\lambda\colon \dot{\mathbf{U}} [\geq \lambda]/ \dot{\mathbf{U}} [> \lambda] \ra {\rm End}(V_\lambda)$} (see~\eqref{pimuiso}). In \cite[Section~29.3]{Lusztig}, Lusz\-tig groups this isomorphism and its properties under the general term of {\it refined Peter--Weyl theorem}. We are going to reinterpret it in terms of $\Oo_A$. First observe that

\begin{lem} The $A$-modules $\Oo_A(\leq \lambda)$ and $\Oo_A(< \lambda)$ are $U_A^{\rm res}$-bimodules, and the surjective map
\begin{gather}\label{dlambdamap}
d_\lambda\colon\ \Oo_A(\leq \lambda)\longrightarrow {\rm Hom}\bigl( \dot{\mathbf{U}}_A [\geq \lambda]/ \dot{\mathbf{U}}_A [> \lambda],A\bigr),\qquad \alpha\longmapsto\langle \alpha , \cdot \ \rangle
\end{gather}
descends to an isomorphism of $U_A^{\rm res}$-bimodules $\bar{d}_\lambda$ on $\Oo_A(\leq \lambda)/\Oo_A(< \lambda)$.
\end{lem}
\begin{proof} For every $\alpha\in \Oo_A(\leq \lambda)$, $x,y\in U_A^{\rm res}$, and $b\in \dot{\mathbf{B}} [\mu]$ with $\mu \nleqslant \lambda$, we have $xby \in \dot{\mathbf{U}}_A [\geq \mu]$. Since $\dot{\mathbf{U}}_A [\geq \mu]= \bigoplus_{\eta\geq \mu}\ A\dot{\mathbf{B}}[\eta]$ and $\eta\geq \mu$ implies $\eta \nleqslant \lambda$, it follows that $\langle xby, \alpha\rangle=0$, i.e., $(x\rhd \alpha \lhd y)(b)=0$. This shows $x\rhd \alpha \lhd y \in \Oo_A(\leq \lambda)$. The same proof applies as well to $\Oo_A(<\lambda)$, whence the first claim. Since $\dot{\mathbf{U}} [\geq \lambda]$ and $\dot{\mathbf{U}} [> \lambda]$ are two-sided ideals of $\dot{\mathbf{U}}$, $\dot{\mathbf{B}}$ is a basis of~$\dot{\mathbf{U}}_A$, and the $A$-modules $\dot{\mathbf{U}}_A [\geq \lambda]$ and $\dot{\mathbf{U}}_A [> \lambda]$ are spanned by \smash{$ \coprod_{\lambda'\geq \lambda} \dot{\mathbf{B}} [\lambda']$} and \smash{$ \coprod_{\lambda'> \lambda} \dot{\mathbf{B}} [\lambda']$}, both are two-sided ideals of $\dot{\mathbf{U}}_A$, and $\dot{\mathbf{U}}_A [\geq \lambda]/ \dot{\mathbf{U}}_A [> \lambda]$ inherits the quotient $U_A^{\rm res}$-bimodule structure. Clearly, the map $d_\lambda$ is well defined, it is a morphism of $U_A^{\rm res}$-bimodules, and its kernel contains $\Oo_A(< \lambda)$. Bijectivity of $\bar{d}_\lambda$ comes by comparing the cardinality of canonical bases: $\Oo_A(\leq \lambda)/\Oo_A(< \lambda)$ has the basis formed by the cosets of the elements of the basis $(\dot{\mathbf{B}}[\lambda])^*$ of~\smash{${}_A \overset{\raisebox{-1.5pt}{\scriptsize$\smallbullet$}}{C} (\lambda')$}, and $\dot{\mathbf{U}}_A [\geq \lambda]/ \dot{\mathbf{U}}_A [> \lambda]$ the basis formed by the cosets of the elements of $\dot{\mathbf{B}}[\lambda]$, all cosets being non-zero and pairwise distinct.
\end{proof}

Since $\dot{\mathbf{U}}_A$ preserves the canonical basis $\underline{\mathbf{B}}_\lambda$ of ${}_AV_\lambda$, $\bar{\pi}_\lambda$ descends to an isomorphism of $U_A^{\rm res}$-bimodules \smash{$\bar{\pi}_\lambda\colon \dot{\mathbf{U}}_A [\geq \lambda]/ \dot{\mathbf{U}}_A [> \lambda] \ra {\rm End}({}_A V_\lambda)$}. We thus get exact sequences of $U_A^{\rm res}$-bi\-mod\-ules%
\[\xymatrix{0 \ar[r] & \dot{\mathbf{U}}_A [> \lambda] \ar[r] & \dot{\mathbf{U}}_A [\geq \lambda] \ar[r]^{\bar{\pi}_\lambda\quad \quad } & {\rm End}({}_A V_\lambda) \longrightarrow 0}\]
and
\begin{equation}\label{exactsequence}
\xymatrix{0 \ar[r] & \Oo_A(< \lambda) \ar[r] & \Oo_A(\leq \lambda) \ar[rr]^{(\bar{\pi}_\lambda^{-1})^*\circ d_\lambda\quad \ } & & \left({\rm End}({}_A V_\lambda)\right)^* \ar[r] & 0.}
\end{equation}
They split as sequences of $A$-modules but not as sequences of bimodules. In fact,
\begin{align}
\left({\rm End}({}_A V_\lambda)\right)^*:={}& {\rm Hom}({\rm End}({}_A V_\lambda),A )\nonumber\\
\cong{} &{\rm Hom}({}_A^{\omega} V_\lambda \bigotimes_A {}_A V_\lambda,A ) = {}_AV_\lambda^{\rm up} \bigotimes_A ({}_A^{\omega}V_\lambda)^{\rm up},\label{identcoeffmatup}
\end{align}
 with the ``\ ${}^{\rm up}$\ '' structure defined in \eqref{upmodule}, and corresponding basis $\underline{\mathbf{B}}_\lambda^{\rm up} \otimes ({}^{\omega}\underline{\mathbf{B}}_\lambda)^{\rm up}$. Moreover, the exact sequence \eqref{exactsequence} shows that this $A$-module of matrix coefficients, regarded as an $A$-submodule of $\Oo_A$ by means of the coefficient map $\Phi := \bigoplus_{\lambda\in P_+} \Phi_\lambda$ (see \eqref{pairingOAKashi}), is contained in $\Oo_A(\leq \lambda)$. This for all $\lambda'\leq \lambda$ yields $ \bigoplus_{\lambda'\leq \lambda} \left({\rm End}({}_A V_{\lambda'})\right)^* \subset \Oo_A(\leq \lambda)$. Now, using the isomorphism~$\bar{\pi}_\lambda$, we get
%\begin{equation}\label{rankOAleq}
\[{\rm rank}_A(\Oo_A(\leq \lambda)) = \sum_{\lambda'\leq \lambda} {\rm Card}\bigl( \dot{\mathbf{B}}[\lambda']\bigr) = \sum_{\lambda'\leq \lambda} {\rm rank}({}_AV_{\lambda'})^2\]
%\end{equation}
and therefore
\begin{equation}\label{rankOqleq}
{\rm dim}_{\mc(q)}\bigl(\Oo_A(\leq \lambda) \bigotimes_A \mc(q)\bigr) = \sum_{\lambda'\leq \lambda} {\rm dim}(V_{\lambda'})^2 = \sum_{\lambda'\leq \lambda} {\rm dim}((C(\lambda')),
\end{equation}
where as usual $C(\lambda')$ denotes the space of matrix coefficients of $V_{\lambda'}$ (see \eqref{defCmu}). It follows
\begin{equation}\label{tensorCOA}
\Oo_A(\leq \lambda) \bigotimes_A {\mathbb C}(q)=\bigoplus_{\lambda'\leq \lambda}C(\lambda'), \qquad\Oo_A(<\lambda) \bigotimes_A {\mathbb C}(q)=\bigoplus_{\lambda'< \lambda}C(\lambda').
\end{equation}
However, in general \smash{${}_A \overset{\raisebox{-1.5pt}{\scriptsize$\smallbullet$}}{C} (\lambda)\bigotimes_A {\mathbb C}(q)$} is not equal to $C(\lambda)$, \smash{${}_A \overset{\raisebox{-1.5pt}{\scriptsize$\smallbullet$}}{C} (\lambda)$} is not an $A$-sublattice of~$C(\lambda)$, and \smash{${}_A \overset{\raisebox{-1.5pt}{\scriptsize$\smallbullet$}}{C} (\lambda)$} is not a $U_A^{\rm res}$-bimodule (it is because of this discrepancy that we have introduced the dot notation ``\ \smash{${}^{\smallbullet}$}\ ''). For instance, we can see the first two facts in the case of $\mathfrak{g} = {\mathfrak{sl}_2}$, by inverting the system of identities \eqref{formulesl2} for all $0\leq l$, $m\leq n$ (or more simply by considering the identity $v_0^* \otimes v_0=[2]_qbc +1$ from \eqref{matrixcoefV2}). For the third fact, we have ${1_2 E\in \dot{\mathbf{B}} [2]}$ (see~\eqref{2sidedcellsl2}), so~${((1_2 E)^* \lhd E)(1_0) \!=\!\langle \Delta((1_2 E)^* ), E\otimes 1_0\rangle\!= \! \langle(1_2 E)^* , E1_0\rangle\!=\!\langle(1_2 E)^*, 1_2E\rangle\!=\!1}$ since $E1_0\! =\! 1_2E$. Therefore, \smash{$(1_2 E)^* \lhd E \notin {}_A \overset{\raisebox{-1.5pt}{\scriptsize$\smallbullet$}}{C} (2)$}.

From the formulas \eqref{formulesl2} and Appendix~\ref{lowupbases}, we can observe the isomorphism~\eqref{identcoeffmatup} in the case of $\mathfrak{g} = {\mathfrak{sl}_2}$. More simply, by projecting the matrix \eqref{matrixcoefV2} onto~$({\rm End}({}_A \!V_2))^*$ the entries are unchanged except the $(1,1)$ entry, which becomes $[2]_q bc$. All factors $[2]_q$ in the middle column disappear if one uses matrix coefficients in the upper canonical basis of~$V_2$, which is $v_0^{\rm up}:=v_0$, $v_1^{\rm up}:=[2]_q^{-1}v_1$, $v_2^{\rm up}:=v_2$ in the notations of \eqref{formulesl2}, since we have
\smash{$v^l(\pi_2(\cdot) \ v_m) = [\delta_{m,1}+1]_q \big\langle v_l^{\rm up}, \cdot v_m^{\rm up}\big\rangle$} for $l,m\in \{0,1,2\}$, where $\langle\ ,\ \rangle$ is the pairing \eqref{proppairinglambda}. Thus, in this particular example of \smash{$({\rm End}({}_A V_2))^*$} we see explicitly the identification of the basis $(\bar{\pi}_2^*)^{-1} \circ d_2(\dot{\mathbf{B}} [2]^*)$ and~${\underline{\mathbf{B}}_2^{\rm up} \otimes ({}^{\omega}\underline{\mathbf{B}}_2)^{\rm up}}$.

Summing up this discussion, the Lusztig refined Peter--Weyl theorem of \cite[Section~29.3]{Lusztig}, implies the following.

\begin{teo} As an $A$-module we have a direct sum decomposition
\begin{equation}\label{OAsurA} \Oo_A = \bigoplus_{\lambda\in P_+} {}_A \overset{\raisebox{-1.5pt}{\scriptsize$\smallbullet$}}{C} (\lambda),
\end{equation}
as $U_A^{\rm res}$-bimodules we have a $($directed by inclusion, and non direct$)$ sum
\begin{equation}\label{OAbimod}
\Oo_A = \sum_{\lambda\in P_+} \Oo_A(\leq \lambda),
\end{equation}
and the composition factors of $\Oo_A$ are the bimodules
\begin{equation}\label{OAfacteurs}
\left({\rm End}({}_A V_\lambda)\right)^* \cong \left({}_A^{\omega} V_\lambda \otimes {}_A V_\lambda \right)^*
\end{equation}
for every $\lambda\in P_+$, each of multiplicity $1$. \end{teo}
\begin{Remark}{\rm The above filtration and its composition factors appear in disguised manner as {\it good filtration} in \cite{APW} and \cite{Paradowski} (see also \cite{VV}).}
\end{Remark}
Because $\dot{\mathbf{B}}$ is formed by weight vectors for the left and right action of $U_q^{\rm ad}(\mathfrak{h})$ (see \eqref{Bpointunion}), the same is true of \smash{$\dot{\mathbf{B}}^*$} and \eqref{OAsurA} can thus be refined into a weight space decomposition
\begin{equation}\label{OAweightdecomp} \Oo_A = \bigoplus_{\mu,\nu\in P}\bigoplus_{\lambda\in P_+} \big({}_A \overset{\raisebox{-1.5pt}{\scriptsize$\smallbullet$}}{C} (\lambda)\big)_{\mu,\nu}.
\end{equation}

Now recall the property \eqref{Bpointunion}. Consider in particular the finite subsets $\dot{\mathbf{B}}_{0,\varpi_i}$ and \smash{$\dot{\mathbf{B}}_{\varpi_i,0}$} associated to the fundamental weights $\varpi_i$, $i=1,\dots, m$. The map \smash{$u\mapsto u({}^\omega v_0\otimes v_{w_0(\varpi_i)})$}, \smash{$u\in \dot{\mathbf{U}}$}, allows one to identify $\dot{\mathbf{B}}_{0,\varpi_i}$ with the canonical basis $\underline{\mathbf{B}}_{\varpi_i}$ of ${}^\omega V_{0}\otimes V_{\varpi_i} \cong V_{\varpi_i}$, and therefore with a uniquely determined finite subset $\mathbf{B}_{\varpi_i}$ of the canonical basis $\mathbf{B}^-$ of $U_q^{\rm ad}(\mathfrak{n}_-)$; similarly, one can identify \smash{$\dot{\mathbf{B}}_{\varpi_i,0}$} with a uniquely determined finite subset ${}^\omega \mathbf{B}_{\varpi_i}$ of the canonical basis $\mathbf{B}^+$ of $U_q^{\rm ad}(\mathfrak{n}_+)$. The elements of $\dot{\mathbf{B}}_{0,\varpi_i}$ and $\dot{\mathbf{B}}_{\varpi_i,0}$ are respectively of the form $b^-1_{\varpi_i}$ and $b^+1_{-\varpi_i}$, where $b^- \in \mathbf{B}_{\varpi_i}$ and $b^+\in {}^\omega \mathbf{B}_{\varpi_i}$, and we have (see \cite[Proposition~3.3 and Section~3.4]{Lusztig3}):

\begin{prop}\label{genOA} The algebra $\Oo_A$ is finitely generated. A system of generators is provided by the elements $a^*\in \dot {\mathbf{B}} {}^*$, where $a \in \bigcup_{i=1}^m (\dot{\mathbf{B}}_{0,\varpi_i}\cup \dot{\mathbf{B}}_{\varpi_i,0})$.\end{prop}
Note that the above system of generators of $\Oo_A$ has
$2 \sum_{i=1}^m \dim(V_{\varpi_i})$
 elements. In fact, recall that $\varphi\colon U_q^{\rm ad} \ra U_q^{\rm ad}$ is the anti-automorphism given by $\varphi(E_i) = F_i$, $\varphi(F_i) = E_i$, $\varphi(K_\lambda) = K_{\lambda}$. Denote by $v_{-\varpi_i}$ and $f_{-\varpi_i}$ the canonical lowest-weight vectors of the highest weight modules~${V_{-w_0(\varpi_i)}}$ and $V_{-w_0(\varpi_i)}^r$, respectively, and put the superscript `` ${}^{\rm up}$ '' for the upper canonical basis vectors.
\begin{lem} For every $b^- \in \mathbf{B}_{\varpi_i}$ and $b^+\in {}^\omega \mathbf{B}_{\varpi_i}$, we have
\begin{align}
 &(b^-1_{\varpi_i})^* = \Phi_{\varpi_i}((f_{\varpi_i} \varphi(b^-))^{\rm up} \otimes v_{\varpi_i}), \label{formgenOA1} \\
 &\big(b^+1_{-\varpi_i}\big)^* = \Phi_{-w_0(\varpi_i)}\big(\big(f_{-\varpi_i} \varphi\big(b^+\big)\big)^{\rm up} \otimes v_{-\varpi_i}\big). \label{formgenOA2} \end{align}
In other words, $(b^-1_{\varpi_i})^*$ and $\big(b^+1_{-\varpi_i}\big)^*$ are the matrix coefficients lying on the first and last columns of the matrix representations in the upper canonical bases of the spaces $V_{\varpi_i}$, ${i=1,\dots, m}$.
\end{lem}

\begin{proof} This can be checked by using the isomorphism \eqref{isocrystal}. The key observation is that
 \[\langle \Phi_\lambda(f_\lambda \otimes v_\lambda), 1_\mu\rangle = \langle f_\lambda , 1_\mu v_\lambda \rangle_\lambda = \delta_{\lambda,\mu}\]
 for every $\lambda\in P_+$, $\mu\in P$, and therefore $\Phi_\lambda(f_\lambda \otimes v_\lambda) = 1_\lambda^*$. Then the computation proceeds by using the equivariance of $\Phi$ under the action of \smash{$U_q^{\rm ad} \otimes (U_q^{\rm ad}){}^{\rm op}$}, the fact that $\langle\cdot,\cdot \rangle$ dualizes the bimodules structures on $\Oq$ and $\dot{\mathbf{U}}$, and the description of the associated Kashiwara operators on $\mathcal{B}(\Oo_q)$ and \smash{$\mathcal{B}\bigl(\dot{\mathbf{U}}\bigr)$}. Here is an alternative argument. By the very definition of the sets~\smash{$\dot{\mathbf{B}} [\lambda]$} we have \smash{$b^-1_{\varpi_i}\in \dot{\mathbf{B}} [\varpi_i]$}, \smash{$b^+1_{-\varpi_i}\in \dot{\mathbf{B}} [-\omega_0(\varpi_i)]$}. We wish to check if their duals~$(b^-1_{\varpi_i})^*$,~$(b^+1_{-\varpi_i})^*$ coincide with the elements of $\Oo_A$ on the right sides of \eqref{formgenOA1} and \eqref{formgenOA2}. As already noticed after~\eqref{exactsequence}, by the isomorphism $\Oo_A(\leq \lambda)/\Oo_A(<\lambda) \cong {\rm End}({}_AV_\lambda)^*$ every matrix coefficient of~${}_AV_\lambda$ belongs to $\Oo_A(\leq \lambda)$. Now, the $A$-modules $\Oo_A(\leq \varpi_i)$ and $\Oo_A(\leq -w_0(\varpi_i))$ are generated by~${\dot{\mathbf{B}}[\varpi_i]^*}$ and $\dot{\mathbf{B}} [-w_0(\varpi_i)]^*$, respectively. Because $((\bar{\pi}_{\lambda}^*)^{-1} \circ d_{\lambda})(\dot{\mathbf{B}} [\lambda]^*)$ coincides with~${\underline{\mathbf{B}}_{\lambda}^{\rm up} \otimes ({}^{\omega}\underline{\mathbf{B}}_{\lambda})^{\rm up}}$, the conclusion follows.
 \end{proof}

Note that the same argument implies that, for every $\lambda\in P_+$, any matrix coefficient of $V_\lambda$ in the upper canonical basis and vanishing on the elements of $\dot{\mathbf{B}} [\lambda']$ for $\lambda'<\lambda$ must belong to~$\dot{\mathbf{B}} [\lambda]^*$. For instance, in the ${\mathfrak{sl}_2}$ case, $\Oo_A(\leq 2)$ has canonical basis $\dot{\mathbf{B}} [0]^* \coprod \dot{\mathbf{B}} [2]^*$, so the matrix coefficients of $V_2$ vanishing on $1_0$ belong to $\dot{\mathbf{B}} [2]^*$. This can be observed in \eqref{matrixcoefV2}, using the comments in the paragraph before \eqref{OAsurA}.


Though the $A$-module ${}_AV_\mu \bigotimes_A {}_AV_\nu$ has no decomposition like \eqref{decomprep}, we can refine the map $C(\mu) \otimes C(\nu) \ra C(\mu+\nu)$ in \eqref{projOq} to an $A$-linear map defined on \smash{${}_A \overset{\raisebox{-1.5pt}{\scriptsize$\smallbullet$}}{C} (\mu) \bigotimes_A {}_A \overset{\raisebox{-1.5pt}{\scriptsize$\smallbullet$}}{C} (\nu)$}. Indeed, there is a unique injective morphism of $U_A^{\rm res}$-modules $\mathfrak{T}_{\mu,\nu}\colon{}_AV_{\mu+\nu}\ra {}_AV_\mu \bigotimes_A {}_AV_\nu$, which is given by $\mathfrak{T}_{\mu,\nu}(v_{\mu+\nu}) = v_{\mu} \otimes v_{\nu}$ \cite[Proposition 25.1.2\,(a)--(b)]{Lusztig}. It defines a morphism of based modules
\[(V_{\mu+\nu},\underline{\mathbf{B}}_{\mu+\nu}) \ra (V_\mu \otimes V_\nu ,\underline{\mathbf{B}}_{\mu} \,\lozenge\, \underline{\mathbf{B}}_\nu),
\]
where $\underline{\mathbf{B}}_{\mu} \,\lozenge\, \underline{\mathbf{B}}_\nu := \{b\,\lozenge\, b', b\in \underline{\mathbf{B}}_{\mu}, b'\in \underline{\mathbf{B}}_{\nu} \}$ \cite[Proposition~27.1.7]{Lusztig}. Hence, $\mathfrak{T}_{\mu,\nu}$ is a~split $A$-li\-near map, i.e., there exists a $A$-linear map $\mathfrak{S}_{\mu,\nu}\colon {}_AV_\mu \bigotimes_A {}_AV_\nu \ra {}_AV_{\mu+\nu} $ such that ${\mathfrak{S}_{\mu,\nu}\circ \mathfrak{T}_{\mu,\nu} ={\rm id}}$. Note that $\mathfrak{S}_{\mu,\nu}$ is not a $U_A^{\rm res}$-morphism. Similarly, the unique morphism of $U_A^{\rm res}$-modules \smash{${}^\omega\mathfrak{T}_{\mu,\nu}\colon{}_A^\omega V_{\mu+\nu}\ra {}_A^\omega V_\mu \bigotimes_A {}_A^\omega V_\nu$} is a split injection. Define \smash{$\rho_{\mu',\mu''}\colon \dot{\mathbf{U}}_A \ra {}_A^\omega V_{\mu'} \bigotimes_A {}_AV_{\mu''}$} by
\[
\rho_{\mu',\mu''}(u)= u\biggl({}^\omega v_{w_0(\mu')}\bigotimes_A v_{w_0(\mu'')}\biggr),\]
and $\rho_{\mu',\mu'',\nu',\nu''}\colon \dot{\mathbf{U}}_A{}^{ \hat{\otimes}2} \ra {}_A^\omega V_{\mu'} \bigotimes_A {}_AV_{\mu''} \bigotimes_A {}_A^\omega V_{\nu'} \bigotimes_A {}_AV_{\nu''}$ by
\[\rho_{\mu',\mu'',\nu',\nu''}(u)= u\biggl({}^\omega v_{w_0(\mu')}\bigotimes_A v_{w_0(\mu'')}\bigotimes_A {}^\omega v_{w_0(\nu')}\bigotimes_A v_{w_0(\nu'')}\biggr).\]
Define $\tau_{\mu',\mu'',\nu',\nu''}\colon {}_A^\omega V_{\mu'+\nu'} \bigotimes_A {}_AV_{\mu''+\nu''} \ra {}_A^\omega V_{\mu'} \bigotimes_A {}_AV_{\mu''} \bigotimes_A {}_A^\omega V_{\nu'} \bigotimes_A {}_AV_{\nu''}$ by
\[\tau_{\mu',\mu'',\nu',\nu''} = \big(1\otimes \hat{R}^{-1} \otimes 1\big)({}^\omega \mathfrak{T}_{\mu',\nu'} \otimes \mathfrak{T}_{\mu'',\nu''}).\]
It is an injective morphism of $U_A^{\rm res}$-modules. In \cite[Section 1.13]{Lusztig3}, Lusztig proved that $\tau_{\mu',\mu'',\nu',\nu''}$ is a split $A$-linear map (\cite{Lusztig3} uses $\hat{R}$ instead of $\hat{R}^{-1}$, since our coproducts on $U_q^{\rm ad}$ are opposite), and that it satisfies{\samepage
\begin{equation}\label{commutdiagLusztig}
\tau_{\mu',\mu'',\nu',\nu''} \rho_{\mu'+\mu'',\nu'+\nu''} = \rho_{\mu',\mu'',\nu',\nu''}\Delta,
\end{equation}
where $\Delta$ is the coproduct of $\dot{\mathbf{U}}_A$, see \eqref{defcoprod}.}

Now take $\mu := \mu'=\mu''$, $\nu:= \nu'=\nu'' \in P_+$, and put $\tau_{\mu,\nu} := \tau_{\mu,\mu,\nu,\nu}$. It follows from the classical decomposition~\eqref{decomprep} over $\mc(q)$, and \eqref{projOq} and \eqref{tensorCOA}, that the product of $\Oo_A$ yields a~map $m\colon \Oo_A(\leq \mu) \bigotimes_A \Oo_A(\leq \nu) \ra \Oo_A(\leq \mu+\nu)$.

Denote the projection map \smash{$p_{\mu+\nu}\colon \Oo_A(\leq \mu+\nu) \ra {}_A \overset{\raisebox{-1.5pt}{\scriptsize$\smallbullet$}}{C} (\mu+\nu)$}, define \smash{${}_A \overset{\smallbullet}{\tau}_{\mu,\nu}:= p_{\mu+\nu} \circ m$}, and put
\[\xymatrix{\pi_\lambda'\colon\ \Oo_A(\leq \lambda)\ar[r] & \Oo_A(\leq \lambda)/\Oo_A(< \lambda)\ar[rr]^{\quad (\bar{\pi}_\lambda^*)^{-1}\circ \bar{d}_\lambda} & & ({\rm End}({}_A V_\lambda) )^*},\]
where the first map is the quotient map. Consider the diagram
\[\xymatrix{ {}_A \overset{\raisebox{-1.5pt}{\scriptsize$\smallbullet$}}{C} (\mu) \otimes {}_A \overset{\raisebox{-1.5pt}{\scriptsize$\smallbullet$}}{C} (\nu) \ar[d]^{ \pi_\mu' \otimes \pi_\nu'} \ar[r]^{{}_A \overset{\smallbullet}{\tau}_{\mu,\nu}} & {}_A \overset{\raisebox{-1.5pt}{\scriptsize$\smallbullet$}}{C} (\mu+\nu) \ar[d]^{\pi_{\mu+\nu}'} \\ ({\rm End}({}_A V_\mu) )^* \otimes ({\rm End}({}_A V_\nu) )^* \ar[r]^{\hspace*{1cm} \tau_{\mu,\nu}^t} & ({\rm End}({}_A V_{\mu +\nu}) )^*,}\]
where $\tau_{\mu,\nu}^t$ is the transpose of Lusztig's map $\tau_{\mu,\nu}$.
\begin{prop} \label{teoLuzstigOAscinde} The map \smash{${}_A \overset{\smallbullet}{\tau}_{\mu,\nu} \colon {}_A \overset{\raisebox{-1.5pt}{\scriptsize$\smallbullet$}}{C} (\mu) \bigotimes_A {}_A \overset{\raisebox{-1.5pt}{\scriptsize$\smallbullet$}}{C} (\nu)\ra {}_A \overset{\raisebox{-1.5pt}{\scriptsize$\smallbullet$}}{C} (\mu+\nu)$} is split as an $A$-linear map and the above diagram is commutative.
\end{prop}
\begin{proof} The commutativity of the diagram comes from equation~\eqref{commutdiagLusztig}. The epimorphism $\pi_\lambda'$ is injective on \smash{${}_A \overset{\raisebox{-1.5pt}{\scriptsize$\smallbullet$}}{C} (\lambda)$}, and maps the canonical basis elements to the elements of the upper canonical basis $\underline{\mathbf{B}}_\lambda^{\rm up} \otimes ({}^{\omega}\underline{\mathbf{B}}_\lambda)^{\rm up}$. By Lusztig's results recalled above, the epimorphism $\tau_{\mu,\nu}^t$ splits as an $A$-linear map. Therefore, the same is true of \smash{${}_A \overset{\smallbullet}{\tau}_{\mu,\nu}$}.
\end{proof}

We stress that \smash{${}_A \overset{\smallbullet}{\tau}_{\mu,\nu}$} plays for $\Oo_A$ the same role as the map \eqref{projOq} for $\Oo_q$.

Finally, we consider for any $n\geq 1$ the invariant elements of $\Oo_A^{\otimes n}$ endowed with the action~${\rm coad}^r_n$ of $U_A^{\rm res}$, see~\eqref{actionprod} \big(recall that $\Ll_{0,n}=\Oo_q^{\otimes n}$ as $U_q^{\rm ad}$-module\big).

First note that, by definition, $\Oo_A(G^n)$ is the restricted dual of the Hopf algebra $U_A^{\rm res}\big(\mathfrak{g}^{\oplus n}\big)$, associated to its category of type $1$ modules. By ordering the summands of $\mathfrak{g}^{\oplus n}$ we get an isomorphism $U_A^{\rm res}\big(\mathfrak{g}^{\oplus n}\big)\cong U_A^{\rm res}(\mathfrak{g})^{\otimes n}$, and any type $1$ simple $U_A^{\rm res}(\mathfrak{g})^{\otimes n}$-module is isomorphic to ${V_{[\lambda]} := \bigotimes_{i=1}^n V_{\lambda_i}}$ endowed with the componentwise action, for some ${[\lambda] :=(\lambda_1,\dots, \lambda_n)\in P_+^n}$ (this is a~classical fact; see, e.g., \cite[Theorem~3.10.2]{Et}). Therefore, we have an isomorphism ${\Oo_A(G^n)\cong \Oo_A^{\otimes n}}$.
With the same notation $[\lambda]:=(\lambda_1,\dots, \lambda_n)\in P_+^n$, let us put
\begin{align*}%\label{AClambdan}
&{}_A \overset{\raisebox{-1.5pt}{\scriptsize$\smallbullet$}}{C} ([\lambda]) := \bigotimes_{i=1}^n {}_A \overset{\raisebox{-1.5pt}{\scriptsize$\smallbullet$}}{C} (\lambda_i) = \bigoplus_{b\in \bigotimes_{i=1}^n \dot{\mathbf{B}}[\lambda_i]^*} Ab,\\
&\Oo_A(\leq [\lambda]) := \bigotimes_{i=1}^n \Oo_A(\leq \lambda_i)=\bigoplus_{[\lambda'] \in P_+^n, \lambda_i'\leq \lambda_i} {}_A \overset{\raisebox{-1.5pt}{\scriptsize$\smallbullet$}}{C} ([\lambda']).
\end{align*}
We thus obtain a decomposition into based $(U_A^{\rm res}\otimes (U_A^{\rm res})^{\rm op})^{\otimes n}$-modules
\[%\label{decompositionOA}
\Oo_A^{\otimes n} = \sum_{[\lambda] \in P_+^n} \Oo_A(\leq [\lambda]).
\]
Now ${\rm coad}_n^r = ({\rm coad}^r)^{\otimes n}\circ \Delta^{(n-1)}$ gives structures of $U_A^{\rm res}$-modules to $\Oo_A^{\otimes n}$ and $\Oo_A(\leq [\lambda])$. In order to make it a based module, we give it the ``$\lozenge$'' product of the canonical bases of the factors~${\Oo_A(\leq \lambda_i)}$, i.e.,
\[\dot{\mathbf{B}} [[\lambda]]^* := \lozenge_{i=1}^n \bigg(\coprod_{\lambda_i'\leq \lambda_i}\dot{\mathbf{B}} [\lambda_i']^*\bigg).\]
We thus obtain a decomposition into based $U_A^{\rm res}$-modules
\begin{equation}\label{decompositionOAcoad}
\Oo_A^{\otimes n} = \sum_{[\lambda] \in P_+^n} (\Oo_A(\leq [\lambda]), \dot{\mathbf{B}} [[\lambda]]^*),
\end{equation}
with composition factors $ \bigotimes_{i=1}^n \left({\rm End}({}_A V_{\lambda_i})\right)^*$. By the properties of ``$\lozenge$'' products of bases of based modules, the underlying $A$-module is
\begin{equation}\label{decompositionOA2}
\Oo_A^{\otimes n} = \bigoplus_{[\lambda] \in P_+^n} {}_A \overset{\raisebox{-1.5pt}{\scriptsize$\smallbullet$}}{C} ([\lambda]).
\end{equation}
Finally, we state the last property of based modules we need. Let $(M,B)$ be a based module. Recall the notations introduced around \eqref{decbaseisospace}. It is proved in \cite[Proposition 27.1.8]{Lusztig} that for every $\lambda\in P_+$ the submodule $M[\geq\lambda]$ is a sub-based module of~$M$, and that it has the basis
\begin{equation}\label{partbasemodule}
B\cap M[\geq\lambda] = \bigcup_{\lambda'\geq \lambda} B[\lambda'].
\end{equation}
Consider $M[\ne 0] := \bigoplus_{\lambda\ne 0} M[\lambda]$, the largest proper submodule of $M$ that contains no non-zero invariant element. Recall that the space of {\it coinvariants of} $M$ is
\begin{align*} M_{U_q^{\rm ad}} & = M/ M[\ne 0] = M/\mc(q)\big\{um-\varepsilon(u)m, \,m\in M, \,u\in U_q^{\rm ad}\big\}
\end{align*}
that is, the largest quotient of $M$ with trivial action, where $\varepsilon\colon U_q^{\rm ad}\ra \mc(q)$ is the counit. It follows from \eqref{partbasemodule} that $M[\ne 0]$ is a sub-based module of $M$, with the basis $\bigcup_{\lambda\ne 0} B[\lambda]$, and we have (this is, \cite[Proposition~27.2.6]{Lusztig}):
\begin{prop}\label{Linvbased} The quotient map $\pi\colon M \ra M_{U_q^{\rm ad}}$ is a morphism of based modules, where $M_{U_q^{\rm ad}}$ is endowed with the basis $B_{U_q^{\rm ad}} := \pi(B[0])$.
\end{prop}
Keeping the same notations, let ${}_AM \subset M$ be the $A$-module generated by $B$, and let ${{}_AM^* \subset M^*}$ be the $A$-module generated by $B^*$. They are $U_A^{\rm res}$-modules. Denote by $({}_A M^*)^{U_A^{\rm res}}$ the submodule of $U_A^{\rm res}$-invariant elements of ${}_A M^*$, regarded as a right module in the natural way.
\begin{lem} We have a direct sum decomposition of $A$-modules
\begin{equation}\label{decompAcoinv}
{}_A M^* = ({}_A M^*)^{U_A^{\rm res}} \bigoplus_A {}_A N,
\end{equation}
where ${}_A N\subset {}_A M^*$ is the $A$-submodule generated by $\bigcup_{\lambda\ne 0} B[\lambda]^*$.
\end{lem}
\begin{proof} By Proposition~\ref{Linvbased}, the transpose map \smash{$\pi^t\colon (M_{U_q^{\rm ad}})^* \ra M^*$} is a monomorphism mapping the dual basis \smash{$B_{U_q^{\rm ad}}^*$} to the subset $B[0]^*$ of $B^*$. The image of $\pi^t$ is \smash{$(M^*)^{U_q^{\rm ad}}$}. If we set \smash{${}_A M_{U_A^{\rm res}} = \pi({}_A M)$}, then \smash{$\pi^t(({}_A M_{U_A^{\rm res}})^*) = ({}_A M^*)^{U_A^{\rm res}}$} is generated by $B[0]^*$, which concludes the proof.
\end{proof}

Note that, since $B[0]$ is in general not invariant under the action of $U_A^{\rm res}$, ${}_A N$ need not be stable under this action.

We are now ready to draw consequences of this discussion and the previous results. As usual denote by \smash{$\big(\Oo_A^{\otimes n}\big)^{U_A^{\rm res}}$} the subspace of invariant elements of $\Oo_A^{\otimes n}$ for the action ${\rm coad}_n^r$. In the case $n=1$, it is just the center $\mathcal{Z}(\Oo_A)$.

\begin{teo}\label{teoLuzstigOAinv} \smash{$\big(\Oo_A^{\otimes n}\big)^{U_A^{\rm res}}$} is a direct summand of the $A$-module $\Oo_A^{\otimes n}$ for any $n\geq 1$.
\end{teo}

\begin{proof} By equation \eqref{decompositionOAcoad}, it is enough to show that for every $[\lambda]\in P_+^n$ the invariant elements of $\Oo_A(\leq [\lambda])$ form a direct summand, and these summands are compatible with non-empty intersections $\Oo_A(\leq [\lambda]) \cap \Oo_A(\leq [\lambda'])$. Using that $\Oo_A(G^n) \cong \Oo_A^{\otimes n}$ and viewing $P_+^n$ as the weight lattice of $G^n$, it is enough to prove these claims for $n=1$. Given $\lambda\in P_+$ put
\[P_{\lambda} = \{\lambda'\in P_+, \, \lambda' \nleqslant \lambda\},\]
and denote by $\dot{\mathbf{U}}_A [P_\lambda]$ the $A$-submodule of $\dot{\mathbf{U}}_A$ generated by $\coprod_{\lambda' \in P_\lambda} \dot{\mathbf{B}} [\lambda']$. Also, let us put \smash{$\dot{\mathbf{U}} [P_\lambda]= \dot{\mathbf{U}}_A [P_\lambda]\bigotimes_A \mc(q)$}.
The complement $P_+\setminus P_{\lambda}$ is finite, and if $\lambda'\in P_{\lambda}$ and $\lambda''\geq \lambda'$, then $\lambda''\in P_{\lambda}$. By the results of \cite[Section~29.2]{Lusztig}, $ \dot{\mathbf{U}} [P_\lambda]$ is a two-sided ideal, and the quotient algebra $\dot{\mathbf{U}} / \dot{\mathbf{U}} [P_\lambda]$ is finite-dimensional with unit the coset of \smash{$ \sum_{\lambda'\leq \lambda}1_{\lambda'}$}, and it is semisimple, isomorphic to \smash{$\bigoplus_{\lambda'\leq \lambda} {\rm End}(V_{\lambda'})$} (whereas $\dot{\mathbf{U}}_A / \dot{\mathbf{U}}_A [P_\lambda]$ has indecomposable modules, see Example~\ref{V2V20}).
It inherits from $\dot{\mathbf{U}}$ a canonical basis, formed by the non-zero cosets of elements of $\dot{\mathbf{B}}$, and with this basis $\dot{\mathbf{U}} / \dot{\mathbf{U}} [P_\lambda]$ is a based module for the right adjoint action~${\rm ad}^r$. Similarly as for~\eqref{dlambdamap}, we have a morphism of $U_A^{\rm res}$-modules
\[
\tilde{d}_\lambda\colon\ \Oo_A(\leq \lambda)\longrightarrow {\rm Hom} \big( \dot{\mathbf{U}}_A / \dot{\mathbf{U}}_A [P_\lambda],A\big),\qquad \alpha\longmapsto\langle \alpha , \cdot \ \rangle,\]
which is an isomorphism by \eqref{rankOqleq} and the computation \smash{$ {\rm dim}\bigl(\dot{\mathbf{U}} / \dot{\mathbf{U}} [P_\lambda]\bigr) = \sum_{\lambda'\leq \lambda} {\rm dim}(V_{\lambda'})^2$} in~\cite[Section~29.2]{Lusztig}. Applying Proposition \ref{Linvbased} and \eqref{decompAcoinv} to the based module ${M=\dot{\mathbf{U}} / \dot{\mathbf{U}} [P_\lambda]}$, we obtain that the invariant elements of $\Oo_A(\leq \lambda)$ form a direct summand. Finally, for any ${\lambda,\lambda'\!\in\! P_+}$ we have \smash{$\Oo_A(\leq \! \lambda) \cap \Oo_A(\leq \! \lambda')\cong {\rm Hom} \bigl( \dot{\mathbf{U}}_A \!/\bigl(\dot{\mathbf{U}}_A [P_\lambda]+ \dot{\mathbf{U}}_A [P_{\lambda'}]\bigr),A\bigr)$}.
 Applying Proposition \ref{Linvbased} and \eqref{decompAcoinv} to the based module $M:=\dot{\mathbf{U}} /(\dot{\mathbf{U}} [P_\lambda]+ \dot{\mathbf{U}} [P_{\lambda'}])$, we obtain that the invariant elements $({}_A M^*)^{U_A^{\rm res}}$ of $\Oo_A(\leq \lambda) \cap \Oo_A(\leq \lambda')$ form a direct $A$-summand. Since the latter is a~based $U_A^{\rm res}$-submodule of $ \Oo_A(\leq \lambda)$ and $\Oo_A(\leq \lambda')$, this summand is also a direct $A$-summand of $ \Oo_A(\leq \lambda)^{U_A^{\rm res}}$ and $\Oo_A(\leq \lambda')^{U_A^{\rm res}}$. This shows the $A$-modules $ \Oo_A(\leq \lambda)^{U_A^{\rm res}}$ for all $\lambda\in P_+$ match to form the $A$-summand $(\Oo_A)^{U_A^{\rm res}}$ of $\Oo_A$, and thus concludes the proof.
 \end{proof}

\begin{Remark} Let $(M,B)$, $(M', B')$ be based modules, with tensor product $(M\otimes M',B_{\lozenge})$, and $B_{\lozenge}[0]\subset B_{\lozenge}$ the subset in bijection with the canonical basis of the space of coinvariants~\smash{$(M\otimes M')_{U_q^{\rm ad}}$} (see Proposition \ref{Linvbased}). This subset is described in \cite[Proposition 27.3.8]{Lusztig} in terms of $B$ and $B'$. Since $\dot{\mathbf{U}} / \dot{\mathbf{U}} [P_{\lambda}]$ is semisimple with known summands, and the construction of the ``$\lozenge$'' product of canonical bases is associative, one can recursively compute the subset of the canonical basis of $\bigotimes_{i=1}^n \dot{\mathbf{U}} / \dot{\mathbf{U}} [P_{\lambda_i}]$ (endowed with the action dual to ${\rm coad}_n^r$) which is in bijection with the canonical basis of the space of coinvariants. Therefore, a complete (though highly nontrivial) characterization of the basis of \smash{$\big(\Oo_A^{\otimes n}\big)^{U_A^{\rm res}}$} can be obtained. Examples can be found in \cite[Section~27.3.10]{Lusztig}. In the case $\mathfrak{g}={\mathfrak{sl}_2}$, the canonical basis of the dual space~\smash{${\rm End}\big(V_1^{\otimes n}\big)^*$} has been identified in~\cite{FK} with the canonical basis of the Temperley--Lieb algebra $TL_n(q)$.
\end{Remark}

\begin{exa}\label{V2V20}{\rm The simplest case is already instructive. Namely, consider $V_1$ and $V_2$, the simple $U_q^{\rm ad}({\mathfrak{sl}_2})$-modules of type $1$ and dimension two and three.

On $V_1$, we have the lower canonical basis vectors $v_+$ and $v_-$, such that $Kv_+ = qv_+$, $Ev_+=0$, $v_- = Fv_+$. The canonical lower and upper bases of $V_1$ are both $\{v_+,v_-\}$. Using the relation~\eqref{udef}, we see that the elements of $\dot{\mathbf{B}}_{0,1}$ and $\dot{\mathbf{B}}_{1,0}$ are $1_1$, $F1_1$ and $1_{-1}$, $E1_{-1}$, respectively; the dual linear forms generate $\Oo_A({\rm SL}_2)$, they are the matrix coefficients $a$, $c$, $d$ and $b$ respectively. By \eqref{2sidedcellsl2}, we have $\dot{\mathbf{B}} [1] = \dot{\mathbf{B}}_{0,1}\coprod \dot{\mathbf{B}}_{1,0}$.

Next consider $V_2$. On $V_2$, we have the canonical highest weight vector $v_0$ of weight $2$, and lower canonical basis $\underline{\mathbf{B}}_{2}=\{v_0,v_1,v_2\}$, where $v_1=Fv_0$ and $v_2=F^{(2)}v_0$. We have ${\underline{\mathbf{B}}_{2}^{\rm up}=\big\{v_0,[2]_q^{-1}v_1,v_2\big\}}$ (see Appendix~\ref{lowupbases}). We can identify the ambient space of the right module~$V_2^r$ with that of~$V_2$; its highest weight vector is then~$v_0$, and its canonical lower and upper bases are $\underline{\mathbf{B}}_{2}^r=\{v_0,v_1,v_2\}$ and $\underline{\mathbf{B}}_{2}^r{}^{\rm up}=\big\{v_0,[2]_q^{-1}v_1,v_2\big\}$.

Consider now the module ${}^\omega V_1\otimes V_1$. We have
\[
\hat{R} = \sum_{n=0}^\infty \frac{\bigl(q-q^{-1}\bigr)^n}{[n]_q !} q^{n(n-1)/2} E^n\otimes F^n,
\]
 so the matrix of the involution $\Psi=\hat{R}^{-1}\circ \bar{}\ $ in the basis $v_+\otimes v_+, v_+\otimes v_-,v_-\otimes v_+,v_-\otimes v_-$ is
\[\big(\hat{R}^{-1}\circ \bar{}\ \big)_{{}^\omega V_1 ,V_1}=\begin{pmatrix} 1 &0 & 0 & 0\\0& 1& 0 &0\\0&0&1&0\\ q^{-1}-q &0&0& 1 \end{pmatrix}.\]
Therefore, the canonical basis $\underline{\mathbf{B}}_{1,1}$ is formed by the vectors $v_+ \,\lozenge\, v_+ =v_+\otimes v_+ + q^{-1} v_-\otimes v_-$ and $v_+ \,\lozenge\, v_- =v_+\otimes v_-$, $v_- \,\lozenge\, v_+ =v_-\otimes v_+ $, $v_- \,\lozenge\, v_- =v_-\otimes v_-$. Consider the partition $\underline{\mathbf{B}}_{1,1} = \underline{\mathbf{B}}_{1,1}[2] \cup \underline{\mathbf{B}}_{1,1}[0]$. We have $\underline{\mathbf{B}}_{1,1}[2] = \{v_- \,\lozenge\, v_+,v_+ \,\lozenge\, v_+, ,v_+ \,\lozenge\, v_-\}$, which is a basis of the three-dimensional submodule $W_2$ of $V_1\otimes V_1$. Since $\underline{\mathbf{B}}_{1,1}$ is an $A$-basis of ${}_A^\omega V_1\bigotimes_A {}_AV_1$, it follows that the epimorphism \smash{$\tau_{1,1}^t \colon {}_A\overset{\raisebox{-1.5pt}{\scriptsize$\smallbullet$}}{C}(1) \bigotimes_A {}_A\overset{\raisebox{-1.5pt}{\scriptsize$\smallbullet$}}{C}(1)\ra {}_A\overset{\raisebox{-1.5pt}{\scriptsize$\smallbullet$}}{C}(2)$} splits (see Proposition \ref{teoLuzstigOAscinde}). The vector $v_- \,\lozenge\, v_-$ is cyclic, so $\underline{\mathbf{B}}_{1,1}[0]=\{v_- \,\lozenge\, v_-\}$. By the definitions, we have $v_+ \,\lozenge\, v_+ = (1 \,\lozenge_0 \, 1)_{1,1}$, $v_+ \,\lozenge\, v_- = (1 \,\lozenge_0\, F)_{1,1}$, $v_- \,\lozenge\, v_+ = (F \,\lozenge_0\, 1)_{1,1}$, $v_- \,\lozenge\, v_- = (F \,\lozenge_0\, F)_{1,1}$, so the corresponding elements of \smash{$\dot{\mathbf{B}}_{1,1} \subset \dot{\mathbf{B}}$} are respectively $1_{0}$, $1_{-2}F$, $1_{2}E$, and $F1_{2}E = E1_{-2}F$.

The invariant submodule $W_0$ of ${}^\omega V_1\otimes V_1$ is generated by $v' = v_-\otimes v_- - q^{-1}v_+\otimes v_+$. The~$U_A^{\rm res}$-modules ${}_A^\omega V_1\bigotimes_A {}_AV_1$ and $W_2 \oplus W_0$ are not equal, though they are by extending scalars to $\mc(q)$. Indeed, we have
\[
v_+\otimes v_+= [2]_q^{-1}(qv_+ \,\lozenge\, v_+-v') \notin W_2 \oplus W_0.
\]
 The module of coinvariants is \smash{$({}^\omega V_1\otimes V_1)_{U_q^{\rm ad}} = \mc(q)\{\pi(v_- \otimes v_-)\}$},
where as usual $\pi\colon {}^\omega V_1\otimes V_1 \ra \smash{({}^\omega V_1\otimes V_1)_{U_q^{\rm ad}}}$ is the quotient map. The transpose map \smash{$\pi^t\colon (({}^\omega V_1\otimes V_1)_{U_q^{\rm ad}})^* \ra ({}^\omega V_1\otimes V_1)^*$} sends $(v_- \,\lozenge\, v_-)^*$ to the unique $U_q^{\rm ad}$-invariant linear map
\[
{\rm ev}_1\colon\ {}^\omega V_1\otimes V_1 \ra \mc(q)
\]
 such that ${\rm ev}_1(v_- \otimes v_-)=1$.

Note that, since elements of $\dot{\mathbf{U}}_A [\lambda > 2]$ act trivially on modules with all isotypical components of highest weight $\leq 2$, ${}_A^\omega V_1\bigotimes_A {}_AV_1$ is an indecomposable module over $\dot{\mathbf{U}}_A / \dot{\mathbf{U}}_A [\lambda > 2]$ (that is, \smash{$\dot{\mathbf{U}}_A / \dot{\mathbf{U}}_A [P_2]$} in the notations of Theorem~\ref{teoLuzstigOAinv}).}\end{exa}

\subsubsection[Some consequences on L\_{0,n}\^A and M\_{0,n}\^A]{Some consequences on $\boldsymbol{\Ll_{0,n}^A}$ and $\boldsymbol{\Mm_{0,n}^A}$}

Recall from Section \ref{defintform} the definition of the integral forms $\Ll_{0,n}^A$ and $\Mm_{0,n}^A$.
\begin{prop}\label{L0NAfgfree} $\Ll_{0,n}^A$ and $\Mm_{0,n}^A$ are free $A$-modules, and \smash{$\Mm_{0,n}^A$} is a direct summand of the $A$-module \smash{$\Ll_{0,n}^A$}. Moreover, \smash{$\Ll_{0,n}^A$} is a finitely generated ring.
\end{prop}
\begin{proof} Since \smash{$\Ll_{0,n}^A = \Oo_A^{\otimes n}$} as $U_A^{\rm res}$-modules, by \eqref{decompositionOA2} it has the basis \smash{$\bigcup_{[\lambda]\in P_+^n}\dot{\mathbf{B}} [[\lambda]]^*$}. Therefore, $\Ll_{0,n}^A$ is a free $A$-module. Since $A$ is a principal ideal domain, it follows that \smash{$\Mm_{0,n}^A$} is a~free $A$-submodule \cite[Appendix~2.2]{Lang}. By Theorem \ref{teoLuzstigOAinv}, there is a direct sum decomposition as $A$-module
\begin{equation}\label{decomposition OAn}
\Ll_{0,n}^A = \Mm_{0,n}^A \oplus {}_AN,
\end{equation} and the proof identifies a basis of $\Mm_{0,n}^A$ as a subset of $\bigcup_{[\lambda]\in P_+^n}\dot{\mathbf{B}} [[\lambda]]^*$.

Next, consider the question of finite generation. By the formula \eqref{prodL0nform}, it is enough to verify this for ${\mathcal L}_{0,1}^A$, but ${\mathcal L}_{0,1}^A= \Oo_A$ as an $A$-module, and $\Oo_A$ is finitely generated by the matrix coefficients of the fundamental $U_A^{\rm res}$-modules ${}_AV_{\varpi_k}$, $k\in\{1,\dots,m\}$ (see~\eqref{formgenOA1} and \eqref{formgenOA2}). Any monomials in these generators can be written as a $A$-linear combination of monomials in the same generators but with the product of $\Ll_{0,1}^A$, instead of the product $\star$. This follows from the integrality properties of the $R$-matrix, and the formula inverse to \eqref{mmtilde} (see in \cite[Section 3.3 and the formulas (4.6)--(4.8)]{BR}).
\end{proof}

\begin{Remark}\label{noHaarRoverA}\quad
\begin{itemize}\itemsep=0pt
\item[$(a)$] As noted in \eqref{decompAcoinv}, the $A$-module ${}_AN$ in the decomposition \eqref{decomposition OAn} is in general not a $U_A^{\rm res}$-module. Therefore, the $A$-linear projection map $\Rr_A\colon \Ll_{0,n}^A\ra \Mm_{0,n}^A$ such that ${\rm Ker}(\Rr_A)={}_AN$ is not a Reynolds operator, for it does not satisfy the identity $\Rr_A(\alpha \beta) = \alpha \Rr_A(\beta)$ for all $\alpha\in \Mm_{0,n}^A$, $\beta \in \Ll_{0,n}^A$.

\item[$(b)$] Recall \eqref{formuleRh}. In coherence with (a) above, there is no normalized Haar measure on $\Oo_A$ taking values in $A$. Indeed, by extending scalars over $\mc(q)$ it should otherwise coincide with the Haar measure $h\colon \Oo_q \ra \mc(q)$, but in the notations of Example \ref{V2V20} (see also the comments after \eqref{formulesl2}), since $h(v_0^* \otimes v_0)=0$ we have $h(bc)=-1/\big(q+q^{-1}\big)$, whence $h$ cannot be defined on $\Oo_A$.

\item[$(c)$] The Haar measure yields a well-defined $\mathcal{A}_0$-linear map $h\colon \Ll(\Oo_q) \ra \mathcal{A}_0$ (and analogously $\mathcal{A}_0$-linear and $\mathcal{A}_\infty$-linear maps $h\colon \Ll_{\lozenge}\big(\Oo_q^{\otimes n}\big) \ra \mathcal{A}_0$ and $\bar{h}\colon \bar{\Ll}_{\lozenge}\big(\Oo_q^{\otimes n}\big) \ra \mathcal{A}_\infty$ for any~$n\geq 1$, where \smash{$\big(\Ll_{\lozenge}\big(\Oo_q^{\otimes n}\big),\mathcal{B}[[\lambda]]^*\big)$} is the crystal basis at $q=0$ underlying the based $U_q^{\rm ad}$-module~\eqref{decompositionOAcoad}). Indeed, by \eqref{renormweightA} the lattice $\Ll_{\lambda}^r {}^{\rm up} \bigotimes_{\mathcal{A}_0} \Ll_{\lambda}^{\rm up}$ is generated by the matrix coefficients in the canonical bases of $V_\lambda^r$ and $V_\lambda$. Since the normalisation by powers of $q$ is vacuous on the trivial module $V_0^* \otimes V_0$, and $h$ vanishes on $V_\lambda^* \otimes V_\lambda$ for $\lambda \in P_+\setminus \{0\}$, the claim follows.
 \end{itemize}
\end{Remark}

